\documentclass{article}

\usepackage{arxiv}

\usepackage[utf8]{inputenc}
\usepackage[T1]{fontenc}
\usepackage{hyperref}
\usepackage{url}
\usepackage{booktabs}
\usepackage{amsfonts}
\usepackage{amssymb}
\usepackage{amsmath}
\usepackage{nicefrac}
\usepackage{microtype}
\usepackage{cleveref}
\usepackage{graphicx}
\usepackage[numbers]{natbib}
\usepackage{doi}
\usepackage{tikz}
\usepackage{xcolor}

%% ============================================================
%% Dark mode (off by default; toggle the conditional to enable)
%% ============================================================
\newif\ifdarkmode
\darkmodefalse
\definecolor{darkbg}{HTML}{1E1E1E}
\definecolor{darklink}{HTML}{6CB4EE}
\definecolor{darkcite}{HTML}{8FBC8F}
\ifdarkmode
  \definecolor{lightfg}{HTML}{D4D4D4}
  \pagecolor{darkbg}
  \color{lightfg}
  \hypersetup{
    colorlinks=true,
    linkcolor=darklink,
    citecolor=darkcite,
    urlcolor=darklink,
  }
\else
  \definecolor{lightfg}{HTML}{000000}
  \hypersetup{
    colorlinks=true,
    linkcolor=blue,
    citecolor=blue,
    urlcolor=blue,
  }
\fi

%% ============================================================
%% Reference parameters (derived values are computed by pgfmath)
%% ============================================================
% --- Inputs ---
\pgfmathsetmacro{\refWACC}{0.07}            % Weighted-average cost of capital
\pgfmathsetmacro{\refN}{30}                 % Plant operating lifetime [yr]
\pgfmathsetmacro{\refTc}{6}                 % Construction time [yr]
\pgfmathsetmacro{\refG}{0.02}               % Inflation rate
\pgfmathsetmacro{\refFindirect}{0.20}       % CAS30 indirect cost fraction
\pgfmathsetmacro{\refLandIntensity}{0.25}   % Land intensity [acres/MWe]
\pgfmathsetmacro{\refLandCost}{10000}       % Land cost [$/acre]
\pgfmathsetmacro{\refCbase}{100}            % Levelization example base cost [M$]
% --- Derived ---
\pgfmathsetmacro{\refCRF}{\refWACC*(1+\refWACC)^\refN / ((1+\refWACC)^\refN - 1)}
\pgfmathsetmacro{\refIDCfrac}{((1+\refWACC)^\refTc - 1)/(\refWACC*\refTc) - 1}
\pgfmathsetmacro{\refIDCpct}{\refIDCfrac * 100}
\pgfmathsetmacro{\refAone}{\refCbase * (1+\refG)^\refTc}
\pgfmathsetmacro{\refPV}{\refAone * (1 - ((1+\refG)/(1+\refWACC))^\refN) / (\refWACC - \refG)}
\pgfmathsetmacro{\refLevelized}{\refCRF * \refPV}
\pgfmathsetmacro{\refUnderestimate}{(\refLevelized - \refAone) / \refAone * 100}
\pgfmathsetmacro{\refLandTotal}{\refLandIntensity * \refLandCost / 1000}  % = acres/MWe * $/acre * 1000 MWe / 1e6

\title{A Differentiable Costing Framework for Fusion Power Plants}

\newif\ifuniqueAffiliation
\uniqueAffiliationtrue

\ifuniqueAffiliation
\author{
	Tal Rubin \\
	Astera Institute \\
	\texttt{tal.rubin@astera.org} \\
	\And
	Damien Scott \\
	Astera Institute \\
	\texttt{damien.scott@astera.org} \\
}
\fi

\renewcommand{\shorttitle}{Differentiable Fusion Costing}

\hypersetup{
pdftitle={A Differentiable Costing Framework for Fusion Power Plants},
pdfauthor={Tal Rubin, Mallory Snowden, Reid Westwood, Damien Scott},
pdfkeywords={fusion energy, LCOE, cost estimation, automatic differentiation, JAX},
}

\begin{document}
\maketitle

\begin{abstract}
We present \textsc{1costingfe}, a fusion power plant costing framework that
computes the levelized cost of electricity (LCOE) from the ground up.
The code uses the Code of Accounts Structure (CAS) approach, used in
the ARIES program and later formalized by the Generation~IV Economic
Modeling Working Group, and adopted and developed under ARPA-E funding
by Woodruff and colleagues in the pyFECONS code.
\textsc{1costingfe}'s principal contribution is rebuilding the cost
accounts most sensitive to fuel cycle and confinement family, namely
buildings, magnets, drivers, direct energy converters, isotope
separation, and staffing, from current procurement data, vendor
pricing, and cross-industry benchmarks, in place of the heritage
scaling that dominates these accounts in prior fusion cost models.
Accounts that are largely fuel- and concept-agnostic inherit prior
conventions where appropriate.
Defaults represent Nth-of-a-kind costs after learning-curve effects
are applied; first-of-a-kind treatment is supported.
The framework spans the four candidate fuel cycles (D-T, D-D,
D-${}^3$He, p-${}^{11}$B) and the principal confinement families
(tokamak, stellarator, mirror, field-reversed configuration, laser
and Z-pinch inertial systems, and magneto-inertial concepts).
It is implemented in JAX with reverse-mode automatic differentiation,
exposing exact LCOE gradients to support sensitivity analysis,
Monte Carlo uncertainty propagation, and identification of
technological corridors compatible with a 1~cent/kWh LCOE target.
The economics, physics, and cost-account methodology are described
in turn, and the framework is exercised against two reference designs,
ARC and ARIES-AT, as a calibration cross-check, with divergences
attributed to specific accounts where procurement data departs
from heritage scaling.
\end{abstract}

\keywords{fusion energy \and LCOE \and cost estimation \and automatic differentiation \and JAX}


%% ============================================================
\section{Introduction}
\label{sec:intro}

Fusion energy promises abundant, low-carbon baseload power, but its economic
viability remains an open question. Credible cost estimates are essential for
guiding R\&D investment, comparing confinement concepts, and identifying the
design parameters with the greatest leverage on the levelized cost of
electricity (LCOE).

The most widely used costing methodology for fusion power plants is the
Code of Accounts Structure (CAS) developed by Schulte et~al.\
\citep{schulte1978} at Pacific Northwest Laboratory and later adopted by
the ARIES program \citep{miller1998} and the Generation~IV Economic Modeling
Working Group \citep{emwg2007}. This framework decomposes plant costs into
hierarchical accounts covering pre-construction, direct capital, indirect
services, financial charges, and annual operating costs. The pyFECONS
code \citep{woodruff2026, woodruff2026b} provides a Python implementation of this
methodology.

Power plant costing estimates are challenging even for mature technologies
with extensive historical data, with fission plants exhibiting significant
cost overruns and schedule delays \citep{eashgates2020}. Applying this
framework to fusion introduces two further challenges:
\begin{enumerate}
\item No combination of confinement approach (magnetic, inertial,
  magneto-inertial), system (tokamak, stellarator, laser-driven, etc.)
  and fuel cycle (D-T, D-D, D-${}^3$He, or p-${}^{11}$B) has yet
  demonstrated sustained net electricity generation, and it remains unclear
  which will prove commercially viable.
\item Most proposed fusion plant designs are pre-conceptual, with
  incomplete physics models and unfinalized engineering geometry, making
  the subsystem and material inputs to the cost accounts inherently
  uncertain.
\end{enumerate}

To pursue this search systematically, \textsc{1costingfe} is under active development as a
costing framework built on JAX that computes LCOE from customer-level
requirements across all major confinement families and fuel cycles. The
choice of JAX as the computational backend provides automatic
differentiation - exact partial derivatives of LCOE with respect to every
input parameter in a single backward pass - and vectorized evaluation over
parameter sweeps and Monte Carlo samples. These capabilities transform
sensitivity analysis from a computational bottleneck into a routine
operation, enabling rapid identification of the design levers with the
greatest leverage on cost.
\textsc{1costingfe} is the cost-accounting layer of a broader fusion
techno-economic analysis pipeline; the upstream concept-ingestion,
SysML, and code-generation stages are developed separately in the
\texttt{fusion-tea} repository family.

The field is entering a phase in which fusion power plants are being designed,
not only experiments meant to determine the plasma physics and demonstrate scientific breakeven.
A commercially viable fusion power plant must be designed with its economics in mind, which introduces
additional parameters and constraints that must be added to the optimization loop.
The economic parameters (construction time, economic plant lifetime, capital cost and operating costs,
and the cost of capital, to name a few) become critical components in the calculation, in addition to the
particular plasma physics, nuclear engineering, and traditional engineering parameters.

Most open-source fusion costing frameworks are either
\begin{enumerate}
\item physics-agnostic, accepting user inputs for many, if not all, cost
  accounts, or
\item applicable to a single, relatively mature confinement family and fuel
  cycle (e.g.\ D-T tokamaks), using reduced-order plasma models to
  determine fusion power from the coils and geometry and to size subsystems
  such as neutral beam injectors and blankets.
\end{enumerate}
Individual companies often maintain proprietary models that are unavailable
to the public, as they contain sensitive information about their design and
its maturity.

Certain inputs to the cost model are well constrained regardless of design
choice. The fuel nuclear physics: cross sections, energy per reaction, and
products are immutable. Raw material isotope mix and enrichment ease are
well characterized and unlikely to shift dramatically, though some cost
reduction through learning is expected if a fusion plant design is successful.
Materials science for fusion-specific alloys is expected to lag behind deployment,
owing to the current lack of testing facilities for neutron irradiation and high heat fluxes,
but a learning curve is expected as demand for high-performance materials grows.

On the economics side, the cost of capital is a market condition rather
than a design parameter, and a high cost of capital strongly incentivizes
designs with lower capital costs and faster construction, even at the
expense of higher operating costs.

The \textsc{1cfe} initiative to search for corridors through which fusion energy can reach the
optimistic price target of 1~cent/kWh must balance technological optimism regarding the
potential of medium-to-high Technology Readiness Level (TRL) concepts with healthy skepticism
regarding very early stage physics. As an example, untested direct energy conversion methods are
not specifically considered at this stage, but the potential for their development and deployment is not ruled out.

Data availability for fusion power plant costs is low: no commercial plant
has been built, and the few large-scale projects (ITER, NIF) are research
facilities with non-representative cost structures. In this regime, a model
with too many free parameters risks overfitting to sparse calibration points
and misrepresenting the accuracy of its estimates. The framework is therefore
built bottom-up with deliberate parsimony: each cost account is independently
justified from literature and first principles, parameters are introduced
only when supported by data or physical reasoning, and methodology choices
are documented in per-account documentation.
The computation flows from customer requirements (net electric
power, availability, plant lifetime) through power balance, engineering
sizing, per-account costing, and financial aggregation to a single LCOE
figure, supporting cost overrides at any account level to accommodate
concept-specific engineering detail as it becomes available.

The framework is designed to integrate with other tools in the 1cfe ecosystem,
including \textsc{Fusion-TEA}, a SysML~v2 power plant model that can override
costing parameters for specific implementations, and a backcasting dashboard
for target-driven design exploration.

The remainder of this paper is organized as follows.
\Cref{sec:lcoe} presents the economics module,
\Cref{sec:physics-module} the physics module, and
\Cref{sec:cas} the cost account structure, with detailed treatment of
the accounts for which independent justification studies have been completed.
\Cref{sec:tokamak-model} presents the 0D tokamak model.


%% ============================================================
\section{Economics Module}
\label{sec:lcoe}

The levelized cost of electricity is the constant electricity price at
which the present value of revenue equals the present value of all costs
over the plant lifetime. It aggregates annualized capital and operating
costs into a single metric in \$/MWh:

\begin{equation}
\label{eq:lcoe}
\text{LCOE} = \frac{\text{Annualized cost}}{\text{Annual electricity production}}
\end{equation}

\noindent The annualized cost is the sum of annualized capital cost
$C_{\text{annual}}$ and annualized operating cost $M_{\text{annual}}$
(maintenance, fuel, and manpower), both in M\$/yr. The annual electricity
production is:

\begin{equation}
\label{eq:e-annual}
E_{\text{annual}} = 8760 \times P_{\text{net}} \times N_{\text{mod}} \times A
\quad [\text{MWh/yr}]
\end{equation}

\noindent where $P_{\text{net}}$ is the net electric power per module in MW,
$N_{\text{mod}}$ is the number of reactor modules, $A$ is the plant
availability (capacity factor), and 8760 is the number of hours per year.
The computation of $P_{\text{net}}$ from fusion power and
engineering coefficients is described in the physics module
(\cref{sec:physics-module}).

\Cref{fig:cashflow} illustrates the cash-flow structure. Total capital
expenditure (CAPEX) is disbursed during the construction period, after
which revenue from electricity sales and annual operating expenditure
(OPEX) flow in opposite directions over the plant lifetime.

\begin{figure}[htbp]
\centering
\begin{tikzpicture}[x=0.52cm, y=0.45cm]
  % --- timeline ---
  \draw[thick, lightfg] (0,0) -- (22,0);

  % --- construction and licensing period bracket ---
  \draw[thick, stealth-stealth, color=blue!70!black]
    (0,-5.5) -- node[midway, below, font=\small] {Construction time $T_c$} (8,-5.5);

  % --- plant lifetime bracket ---
  \draw[thick, stealth-stealth, color=blue!70!black]
    (8,-5.5) -- node[midway, below, font=\small] {Plant lifetime $n$ years} (22,-5.5);

  % --- CAPEX arrows (uniform disbursements during construction) ---
  \foreach \x in {0,2,4,6} {
    \draw[thick, -{stealth}, color=lightfg] (\x,0) -- (\x,-3.5);
  }
  \node[font=\small, lightfg] at (3,-4.2) {CAPEX};

  % --- Revenue arrows (upward, during operation, spacing=2, starting 2 from COD) ---
  \foreach \x in {10,12,14,16,18,20,22} {
    \draw[thick, -{stealth}, color=green!50!black] (\x,0) -- (\x,3.5);
  }
  \node[font=\small, color=green!50!black] at (16,4.2) {Revenue};

  % --- OPEX arrows (downward, growing to reflect nominal escalation) ---
  \foreach \x/\h in {10/-1.4, 12/-1.6, 14/-1.8, 16/-2.0, 18/-2.2, 20/-2.5, 22/-2.8} {
    \draw[thick, -{stealth}, color=lightfg] (\x,0) -- (\x,\h);
  }
  \node[font=\small, lightfg] at (16,-3.4) {OPEX};

  % --- time markers ---
  \draw[thin, lightfg!60] (8,-0.3) -- (8,0.3);
  \node[font=\scriptsize, lightfg!60, above] at (8,0.15) {COD};
\end{tikzpicture}
\caption{Schematic cash-flow diagram for a fusion power plant.
  Capital expenditure (CAPEX) is disbursed uniformly during the
  construction period $T_c$, accruing interest until the commercial
  operation date (COD). After COD, revenue from electricity sales and
  annual operating expenditure (OPEX) flow over the plant operating
  lifetime of $n$ years. OPEX grows in nominal terms at the inflation
  rate. The LCOE (\cref{eq:lcoe}) is the constant electricity price at
  which the present value of revenue equals the present value of all
  costs.}
\label{fig:cashflow}
\end{figure}

\subsection{Capital Recovery Factor}

The Capital Recovery Factor (CRF) converts a present-value capital
investment into equal annual payments over the plant lifetime:

\begin{equation}
\label{eq:crf}
\text{CRF}(i, n) = \frac{i\,(1+i)^n}{(1+i)^n - 1}
\end{equation}

\noindent where $i$ is the weighted-average cost of capital (WACC) and
$n$ is the plant operating lifetime in years. At reference conditions
($i = \refWACC$, $n = \pgfmathprintnumber[precision=0]{\refN}$~yr),
CRF $= \pgfmathprintnumber[precision=4]{\refCRF}$.

\subsection{Annualized Capital Cost}
\label{sec:capex-annual}

Capital expenditure is disbursed over the construction period $T_c$
before any revenue is generated (\cref{fig:cashflow}). Assuming the
overnight capital cost $C$ is spent in $T_c$ equal annual installments
of $C/T_c$, each installment accrues compound interest from the time
it is spent until the commercial operation date (COD). The future value
of this series at COD is:

\begin{equation}
\label{eq:fv-capex}
\text{FV} = \frac{C}{T_c} \sum_{k=0}^{T_c - 1} (1+i)^{k}
          = \frac{C}{T_c} \cdot \frac{(1+i)^{T_c} - 1}{i}
\end{equation}

\noindent The interest during construction (IDC) is the excess above
the overnight cost, the total financing charge incurred while the
plant earns no revenue:

\begin{equation}
\label{eq:idc-lcoe}
\text{IDC} = \text{FV} - C
           = C \left[ \frac{(1+i)^{T_c} - 1}{i\,T_c} - 1 \right]
\end{equation}

\noindent The total capital investment at COD is therefore
$C + \text{IDC} = \text{FV}$, and the annualized capital cost is:

\begin{equation}
\label{eq:c-annual}
C_{\text{annual}} = \text{CRF}(i, n) \times (C + \text{IDC})
\end{equation}

\noindent At reference conditions ($i = \refWACC$,
$T_c = \pgfmathprintnumber[precision=0]{\refTc}$~yr), the IDC adds
\pgfmathprintnumber[precision=1]{\refIDCpct}\% to the overnight cost.

\subsection{Levelized Annual Operating Cost}
\label{sec:lac}

Annual operating costs are stated in constant
base-year dollars but escalate at the inflation rate $g$ in nominal
terms over the plant lifetime. The levelized annual cost accounts for
this growth using a three-step procedure.

First, the base-year annual cost $C_0$ is inflated to
first-year-of-operation dollars:

\begin{equation}
\label{eq:inflate}
A_1 = C_0 \, (1 + g)^{T_c}
\end{equation}

\noindent where $T_c$ is the construction time (or total project time
for first-of-a-kind plants, which includes licensing).

Second, the present value of the resulting growing annuity is:

\begin{equation}
\label{eq:growing-annuity}
\text{PV} =
\begin{cases}
\displaystyle A_1 \cdot \frac{1 - \left(\frac{1+g}{1+i}\right)^n}{i - g}
  & \text{if } i \neq g \\[12pt]
\displaystyle A_1 \cdot \frac{n}{1+i}
  & \text{if } i = g
\end{cases}
\end{equation}

Third, the levelized annual cost is obtained by applying the CRF:

\begin{equation}
\label{eq:levelized}
C_{\text{level}} = \text{CRF}(i, n) \times \text{PV}
\end{equation}

\noindent The $i = g$ case follows from L'H\^{o}pital's rule and is
implemented using \texttt{jnp.where} for JAX traceability.
For example, at $C_0 = \pgfmathprintnumber[precision=0]{\refCbase}$~M\$,
$i = \refWACC$, $g = \refG$,
$n = \pgfmathprintnumber[precision=0]{\refN}$~yr,
$T_c = \pgfmathprintnumber[precision=0]{\refTc}$~yr, the levelized
cost is \pgfmathprintnumber[precision=1]{\refLevelized}~M\$/yr
compared to \pgfmathprintnumber[precision=1]{\refAone}~M\$/yr from the
na\"{\i}ve single-year inflation formula, a
\pgfmathprintnumber[precision=0]{\refUnderestimate}\% underestimate
when the growing annuity is neglected.


%% ------------------------------------------------------------
\section{Physics Module}
\label{sec:physics-module}

The denominator of the LCOE, \cref{eq:lcoe}, is the annual saleable
electricity, which depends on the net electric power $P_{\text{net}}$ and plant availability.
The purpose of the physics module is computing $P_{\text{net}}$ from
device and plasma parameters, given a confinement concept and fuel cycle.
A computation of the plant availability would also fit in this, or an Engineering module,
but requires a dedicated model incorporating the specifics of a design (e.g. maintainability, accessibility, component failure modes).

The power balance section of the physics module computes $P_{\text{net}}$
from the fusion power $P_{\text{fus}}$ for the four fusion fuels by tracking
every energy channel from the plasma to the grid: energy split between fusion ash and neutrons,
radiation losses, blanket multiplication, thermal and direct energy conversion, and
recirculating power.

Many concepts will deviate significantly from the defaults - a non-Maxwellian
deuterium plasma may alter the burn rates of the secondary D-D reactions,
or a doped plasma may alter the radiation losses or the plasma optical thickness
\citep{frank2022}. So these and other coefficients are exposed as first-class
parameters amenable to sensitivity analysis via automatic differentiation.

Lumped efficiency parameters are used for characterizing the power flow and thermal cycle.
A more involved power plant thermodynamics model could be implemented in the future,
which would allow the selection of a thermal cycle (Brayton vs. Rankine, etc.) and working
fluid.

The power balance above is concept-agnostic: given $P_{\text{fus}}$
and the engineering coefficients, it produces $P_{\text{net}}$
regardless of confinement scheme. To close the loop and determine
$P_{\text{fus}}$ from reactor geometry and plasma parameters, a
concept-specific physics model is required.
\Cref{sec:tokamak-model} (appendix) describes the zero-dimensional tokamak
equilibrium model included as the first such layer; models for other
confinement families will follow the same pattern.

The power balance described next applies to steady-state confinement
systems (tokamaks, stellarators, mirrors). Pulsed confinement
systems (laser IFE, Z-pinch, and inductive-recovery concepts such as
field-reversed configurations) use a separate per-pulse energy
framework described in \cref{sec:pulsed-balance}.

\subsection{Fusion Power Split}

Operating the power plant requires collecting the power leaving the
plasma and converting it to electricity. The total power exiting the
plasma is the fusion power $P_{\text{fus}}$ plus the external heating
power $P_{\text{in}}$ required to sustain the burn. This power leaves
in three forms: charged particles, neutrons, and photon
radiation. Neutrons deposit their energy volumetrically in the blanket
and shielding; charged particles and photons are absorbed by
plasma-facing surfaces. Most of this energy is converted to electricity
via a thermal cycle, though a fraction of the charged-particle energy
may be recovered directly via direct energy conversion (DEC).

The purpose of this calculation is to determine the power in each of
these channels, which depends on the nuclear physics of the fuel burn.
The neutron energy fraction
$f_n \equiv E_n / Q$, where $Q$ is the total energy released by the
primary reaction and $E_n$ is the portion carried by neutrons.
For fuel cycles with secondary burns (D-D, D-${}^3$He), $f_n$ is
generalised to $\bar{E}_n / \bar{E}_{\text{total}}$ using the
effective averages defined in the subsections below.
\Cref{tab:fusion-reactions} summarizes the primary reaction
for each fuel cycle:

\begin{table}[htbp]
\caption{Primary fusion reaction by fuel cycle.}
\label{tab:fusion-reactions}
\centering
\renewcommand{\arraystretch}{1.3}
\begin{tabular}{lccl}
\toprule
Fuel & $Q$ (MeV) & $f_n$ & Products \\
\midrule
D-T           & 17.6  & 0.80 & ${}^4$He + n \\
D-D           & 3.65  & 0.33 & $\begin{cases} {}^3\text{He} + \text{n} \\ \text{T} + \text{p} \end{cases}$ \\
D-${}^3$He    & 18.3  & 0.0  & ${}^4$He + p \\
p-${}^{11}$B  & 8.7   & 0.0  & 3${}^4$He \\
\bottomrule
\end{tabular}
\end{table}

\subsubsection{Deuterium-tritium (D-T)}
The D-T reaction is the simplest case: each fusion event releases a
fixed $Q = 17.6$~MeV, split by momentum conservation into a 3.5~MeV
alpha particle and a 14.1~MeV neutron ($f_n = 0.80$). The alpha ash
(${}^4$He) is inert and does not undergo further reactions, so the
power split is fully determined by the primary kinematics.

\subsubsection{Deuterium-deuterium (D-D)}
The D-D fuel cycle is the most complex fuel cycle. The primary reaction
branches into two channels with approximately equal probability, producing
either tritium and a proton (branch~1) or ${}^3$He and a neutron (branch~2).
The tritium and ${}^3$He products are themselves fusion fuels that react
with the majority species of deuterium in the plasma. The deuterium-tritium
(D-T) reaction has a much higher cross section than the D-D reaction, so the tritium
from branch~1 undergoes secondary D-T fusion with a burn fraction $f_T$ that
approaches 1 at typical D-D burn temperatures. The D-${}^3$He reaction
has a cross section comparable to D-D, so the ${}^3$He product
from branch~2 undergoes secondary D-${}^3$He fusion with a burn fraction 
$f_{{}^3\text{He}}$ that is typically lower than $f_T$.

The effective energy per D-D event, and the split between charged particles and neutrons,
is determined by the burn fractions of these secondary reactions. 
The burn fractions depend on the plasma conditions and confinement time, and are
currently treated as free parameters in the framework, with default values
informed by typical D-D burn conditions in the literature \citep{huba2016}.
A detailed plasma model could be added to compute these burn fractions from plasma parameters.
The effective energy per D-D event, averaged over the two primary branches, becomes:

\begin{equation}
\label{eq:dd-etotal}
\bar{E}_{\text{total}} = \tfrac{1}{2}(Q_1 + Q_2)
  + \tfrac{1}{2} f_T \, Q_{\text{DT}}
  + \tfrac{1}{2} f_{{}^3\text{He}} \, Q_{\text{D}{}^3\text{He}}
\end{equation}

\noindent where $Q_1 = 4.03$~MeV and $Q_2 = 3.27$~MeV are the primary
branch Q-values. Each secondary reaction consumes one additional
deuteron, so the total deuterons consumed per primary D-D event is
$2 + \tfrac{1}{2} f_T + \tfrac{1}{2} f_{{}^3\text{He}}$.
The neutron energy per primary D-D event is

\begin{equation}
\label{eq:dd-eneutron}
\bar{E}_n = \tfrac{1}{2} E_{n,1} + \tfrac{1}{2} f_T \, E_{n,\text{DT}}
\end{equation}

\noindent where $E_{n,1} = 2.45$~MeV is the neutron from primary branch~2
and $E_{n,\text{DT}} = 14.1$~MeV is the neutron from the secondary D-T
reaction; the secondary D-${}^3$He reaction produces only charged
particles and contributes no neutrons.
At default burn fractions ($f_T = 0.97$,
$f_{{}^3\text{He}} = 0.69$), 2.83 deuterons are consumed per primary
event, the effective total energy rises from
3.65~MeV to approximately 18.5~MeV, and the effective
neutron energy fraction $\bar{E}_n / \bar{E}_{\text{total}}$ increases from
0.33 to approximately 0.44 due to the secondary D-T neutrons.

\subsubsection{\texorpdfstring{Deuterium-helium 3 (D-${}^3$He)}{Deuterium-helium 3 (D-3He)}}
The primary D-${}^3$He reaction produces only charged particles
($Q = 18.3$~MeV). At the temperatures required for D-${}^3$He fusion
($T_i = 60$--$100$~keV), a fraction $f_{\text{DD}}$ of fusion events
are D-D side reactions. The D-D events produce tritium (via
D(d,p)T) that re-burns in D-T fusion at burn fraction
$f_T^*$, introducing 14.1~MeV neutrons into an
otherwise neutron-free system. They also produce ${}^3$He
(via D(d,n)${}^3$He) that re-burns in D-${}^3$He fusion at burn
fraction $f_{{}^3\text{He}}^*$, releasing additional
charged-particle energy and reducing the external ${}^3$He demand
of the plant.

The effective total energy and neutron energy per D-${}^3$He event are:

\begin{equation}
\label{eq:dhe3-etotal}
\bar{E}_{\text{total}} = Q_{\text{D}{}^3\text{He}}
  + f_{\text{DD}} \, \bar{E}_{\text{total,DD}}
\end{equation}

\begin{equation}
\label{eq:dhe3-eneutron}
\bar{E}_n = f_{\text{DD}} \, \bar{E}_{n,\text{DD}}
\end{equation}

\noindent where $\bar{E}_{\text{total,DD}}$ follows the structure
of \cref{eq:dd-etotal} including both secondary burns:
\begin{equation*}
\bar{E}_{\text{total,DD}} = \tfrac{1}{2}(Q_1 + Q_2)
  + \tfrac{1}{2} f_T^* \, Q_{\text{DT}}
  + \tfrac{1}{2} f_{{}^3\text{He}}^*
    \, Q_{\text{D}{}^3\text{He}}
\end{equation*}
and $\bar{E}_{n,\text{DD}}$ follows \cref{eq:dd-eneutron}: only the
D-T secondary burn contributes neutrons, since the secondary
D-${}^3$He reaction is aneutronic. The burn-fraction parameters
$f_T^*$ and $f_{{}^3\text{He}}^*$
are kept separate from the D-D fuel cycle's burn fractions and are
concept-specific. The single-pass burn fraction in a thermal plasma
is set by $f_X = n_D \langle\sigma v\rangle_X \tau_p / (1 + n_D
\langle\sigma v\rangle_X \tau_p)$ where $\tau_p$ is the species
particle confinement time; for a steady-state mirror at $T_i = 70$~keV
with $n_D = 3.3 \times 10^{19}$~m${}^{-3}$ and $\tau_p \approx 30$~s,
$f_T^* \approx 0.5$ and $f_{{}^3\text{He}}^* \approx 0.1$, with the latter
smaller because $\langle\sigma v\rangle_{\text{D}{}^3\text{He}}$ is
approximately 13$\times$ smaller than $\langle\sigma v\rangle_{\text{DT}}$ at this
temperature (see
\texttt{examples/dhe3\_burn\_fractions.py}). Pulsed-FRC architectures
with active inter-shot ${}^3$He recovery reach the
asymptotic limit set by the recovery efficiency: at 99\% recovery
per cycle, $f_{{}^3\text{He}}^* \approx 0.99$. Tritium in such a
plant is exhausted (or held until decay) to keep the design
neutron-averse, so $f_T^* \approx 0$. Default values are therefore
concept-specific: steady-state mirror uses $f_T^* = 0.5$,
$f_{{}^3\text{He}}^* = 0.1$; pulsed FRC uses $f_T^* = 0$,
$f_{{}^3\text{He}}^* = 0.99$.
Per primary D-${}^3$He event, the deuteron
consumption is $1 + f_{\text{DD}} (1 + \tfrac{1}{2} f_T^*
+ \tfrac{1}{2} f_{{}^3\text{He}}^*)$ and the
external ${}^3$He consumption is $(1 - f_{\text{DD}})
- \tfrac{1}{2} f_{\text{DD}} f_{{}^3\text{He}}^*$,
the latter accounting for D-D-bred ${}^3$He that re-burns in place
of an external atom.
For the steady-state mirror concept evaluated at the default
$f_{\text{DD}} = 0.131$ (corresponding to a 50/50 D/${}^3$He mix at
$T_i = 70$~keV with Bosch-Hale cross sections \citep{boschhale1992}),
together with the steady-state default $f_T^* = 0.5$ and
$f_{{}^3\text{He}}^* = 0.1$, the effective neutron energy fraction
$\bar{E}_n / \bar{E}_{\text{total}}$ is approximately 6\%, small but
nonzero, and consequential for shielding, activation, and tritium
breeding requirements. See \texttt{examples/dhe3\_mix\_optimization.py}
for the self-consistent calculation that links $f_{\text{DD}}$ to the
plasma D/${}^3$He density ratio and operating temperature.

\subsubsection{\texorpdfstring{Proton-boron 11 (p-${}^{11}$B)}{Proton-boron 11 (p-11B)}}
The primary p-${}^{11}$B reaction produces three alpha particles
($Q = 8.7$~MeV) with a broad energy spectrum. \citet{dmitriev2009} reports
a primary alpha peaking near 4.5~MeV and two secondary alphas from the
decay of an excited ${}^8$Be intermediate at approximately 2.4~MeV each. These alphas
are not inert ash but can undergo further reactions with the boron fuel.
The dominant side reaction is ${}^{11}$B($\alpha$,n)${}^{14}$N ($Q = 0.16$~MeV).
The Coulomb barrier between $\alpha$ and ${}^{11}$B is 3.2~MeV, which is
below the primary alpha energy but above the secondary alpha energies.
Measured total cross sections range from 20 to 240~mb
over alpha energies of 2--6~MeV, with significant resonance structure
\citep{prior2017}. A secondary side reaction,
${}^{11}$B(p,n)${}^{11}$C ($Q = -2.76$~MeV, threshold 3.02~MeV),
is endothermic and accessible to protons in the high-energy tail at a much lower rate
($10^{-5}$ relative to the primary reaction)
\citep{putvinski2019}.

The framework parameterizes these via $f_{\alpha n}$ (probability per
alpha of undergoing ($\alpha$,n)) and $f_{pn}$ (fraction of primary
events that instead undergo (p,n)). Since each primary event produces
three alphas, the effective energies per primary event are:

\begin{equation}
\label{eq:pb11-etotal}
\bar{E}_{\text{total}} = Q_{\text{pB}} + 3 f_{\alpha n} \, Q_{\alpha n}
  + f_{pn} \, Q_{pn}
\end{equation}

\begin{equation}
\label{eq:pb11-eneutron}
\bar{E}_n = 3 f_{\alpha n} \, E_{n,\alpha} + f_{pn} \, E_{n,p}
\end{equation}

\noindent where $Q_{\alpha n} = +0.16$~MeV and $Q_{pn} = -2.76$~MeV
are the side-reaction Q-values, $E_{n,\alpha} \approx 2.2$~MeV and
$E_{n,p} \approx 0.17$~MeV are the neutron kinetic energies from
center-of-mass kinematics.
Both $f_{\alpha n}$ and $f_{pn}$ default to zero:
\citet{ochs2025} show that viable p-${}^{11}$B reactors must remove
alpha ash on timescales much shorter than the energy confinement time
to avoid excessive bremsstrahlung losses, which drastically reduces the
population of MeV-scale alphas coexisting with the boron fuel and thus
suppresses the dominant ${}^{11}$B($\alpha$,n)${}^{14}$N channel
\citep{hora2021}. At the default values, p-${}^{11}$B is treated as
purely aneutronic.

\subsection{Thermal and Electric Power}

The power deposited in the plasma (charged-particle fusion products
$P_{\text{cp}}$ plus external heating $P_{\text{in}}$) exits as
radiation losses (bremsstrahlung, synchrotron, and impurity line
radiation) and transport power deposited on plasma-facing surfaces.
Bremsstrahlung power \citep{wesson2011}:

\begin{equation}
\label{eq:p-brem}
P_{\text{brem}} = 5.35 \times 10^{-37}\; n_e^2 \, Z_{\text{eff}}
  \sqrt{T_e}\; V \quad [\text{W}]
\end{equation}

\noindent where $n_e$ is the electron density in m$^{-3}$, $T_e$ the
volume-averaged electron temperature in keV, $Z_{\text{eff}}$ the
effective charge, and $V$ the plasma volume in m$^3$.

For synchrotron radiation, the optically thin (Larmor) formula
grossly overestimates the net loss in MFE plasmas because most
emission is reabsorbed. The framework uses the global model of
\citet{albajar2001} with the wall-reflectivity correction of
\citet{fidone2001}:

\begin{equation}
\label{eq:p-sync}
P_{\text{sync}} = 3.84 \times 10^{-8}\;
  C(R_w)\; R\, a^{1.38}\, \kappa^{0.79}\, B^{2.62}\,
  n_{e0,20}^{0.38}\; T_{e0}\,(16 + T_{e0})^{2.61}\;
  K\, G
  \quad [\text{MW}]
\end{equation}

\noindent where $T_{e0}$ and $n_{e0,20}$ are central values (keV
and $10^{20}$~m$^{-3}$), $R$ the major radius, $a$ the minor
radius, $\kappa$ the elongation, and $B$ the toroidal field in T.
The correction factor
$C(R_w) = (1-R_w)^{0.62}\bigl[1 + 0.12\,T_{e0}\,(1-R_w)^{0.41}/
p_{a0}^{0.41}\bigr]^{-1.51}$
captures wall reflectivity $R_w$ and optical thickness via
$p_{a0} = 6.04 \times 10^3\, a\, n_{e0,20}/B$.
The profile shape factor $K(\alpha_n, \alpha_T, \beta_T)$ and
aspect-ratio correction $G(A) = 0.93\,[1 + 0.85\,e^{-0.82\,A}]$
account for profile peaking and toroidal geometry.
Although derived for tokamak geometry, the formula is applied to all
MFE concepts: for stellarators the toroidal geometry is directly
analogous, while for mirrors the effective major radius is set to
$R_{\text{eff}} = L/(2\pi)$ (mapping the cylinder length $L$ to a
torus of equivalent volume) with reduced wall reflectivity to account
for the open ends. IFE and MIF set $P_{\text{rad}} = 0$.

The Albajar model assumes a thermal (Maxwell--J\"uttner) electron
distribution. At the high temperatures required for aneutronic fuels
($T_e \gtrsim 100$~keV), synchrotron drag itself suppresses the
high-energy electron tail, which is responsible for most of the
emission at optically thin harmonics \citep{ochs2024}. This
self-consistent kinetic effect can reduce the net synchrotron loss by
up to an order of magnitude relative to the thermal prediction. For
such regimes the user should supply a radiation power override
informed by kinetic Fokker--Planck modeling rather than rely on the
Albajar formula.

\subsubsection{Impurity line radiation.}
In MFE plasmas, line radiation from partially ionized impurities is
often the dominant radiation loss channel. The framework models this
via a steady-state sputtering--transport chain connecting the
first-wall material to the core impurity concentration and radiative
loss.

The wall material is specified as one of six options: tungsten~(W),
carbon~(C), beryllium~(Be), molybdenum~(Mo), silicon carbide~(SiC),
or liquid lithium~(Li). For each material, a simplified Bohdansky
physical sputtering yield \citep{bohdansky1984, eckstein2007} gives
the number of wall atoms released per incident ion at the
sheath-accelerated energy $E_{\text{ion}} \approx 3\,T_{\text{edge}}$:

\begin{equation}
\label{eq:sputtering}
Y(E) = Q \left(1 - \left(\frac{E_{\text{th}}}{E}\right)^{2/3}\right)
       \left(1 - \frac{E_{\text{th}}}{E}\right)^{2}
\end{equation}

\noindent where $E_{\text{th}}$ is the material-dependent threshold
energy and $Q$ is an effective fit coefficient that absorbs the
Thomas--Fermi nuclear stopping cross-section $S_n(\varepsilon)$. Below
threshold ($E < E_{\text{th}}$), $Y = 0$. The steady-state impurity fraction in the core is then:

\begin{equation}
\label{eq:f-impurity}
f_z = Y(3\,T_{\text{edge}}) \cdot
  \frac{A_{\text{fw}}}{V} \cdot
  \frac{\tau_{\text{imp}}}{\tau_E} \cdot f_{\text{screen}}
\end{equation}

\noindent where $A_{\text{fw}}/V$ is the first-wall area to plasma
volume ratio, $\tau_{\text{imp}}/\tau_E$ is the impurity-to-energy
confinement time ratio (default 3.0), and $f_{\text{screen}}$ is the
scrape-off layer screening factor (default 0.01 for high-$Z$
materials, 0.1 for low-$Z$). Lithium is a special case: as a fully
ionized species above 0.1~keV, it produces negligible line
radiation at reactor core temperatures; the primary benefit of a
liquid Li wall is low sputtering yield from continuous surface renewal
and correspondingly low $f_{\text{screen}}$.

Deliberately seeded impurities (e.g.\ neon, argon) for divertor
radiation or edge cooling can be specified directly via their
concentration $f_z$, bypassing the sputtering model.

The impurity line radiation power uses coronal-equilibrium radiative
loss rate coefficients $L_z(T_e)$
\citep{post1977, mavrin2018}, represented as piecewise
power-law fits for each species:

\begin{equation}
\label{eq:p-line}
P_{\text{line}} = n_e^2 \sum_z f_z \, L_z(T_e) \; V \quad [\text{W}]
\end{equation}

\noindent The summation runs over all impurity species (both
wall-derived and seeded). Cooling curves are tabulated for W, C, Be,
Mo, Si, Li, Ne, and Ar.

All new parameters default to no-op values (null wall material, empty
seeded impurities), preserving backward compatibility: when no wall
material is specified, $P_{\text{line}} = 0$ and the radiation model
reduces to bremsstrahlung plus synchrotron as before.

\subsubsection{Plasma energy balance.}

The total radiated power is
$P_{\text{rad}} = P_{\text{brem}} + P_{\text{sync}} + P_{\text{line}}$,
computed at the user-specified electron temperature $T_e$.
In steady state, the plasma energy balance requires that all power
entering the plasma (charged-particle fusion products plus external
heating) equals all power leaving (radiation plus transport to
plasma-facing surfaces):

\begin{equation}
\label{eq:plasma-balance}
P_{\text{cp}} + P_{\text{in}} = P_{\text{rad}} + P_{\text{transport}}
\end{equation}

\noindent where $P_{\text{cp}}$ is the charged-particle (ash) power and
$P_{\text{in}}$ is the external heating power. The transport power is
therefore $P_{\text{transport}} = P_{\text{cp}} + P_{\text{in}} - P_{\text{rad}}$.

If $P_{\text{rad}}$ exceeds $P_{\text{cp}} + P_{\text{in}}$ at the
specified $T_e$, the user-supplied heating power is insufficient to
sustain the plasma against radiation losses. Rather than capping
$P_{\text{rad}}$ implicitly (which would mask an unphysical operating
point), the framework increases the effective heating power to the
minimum needed to close the energy balance:

\begin{equation}
\label{eq:p-input-eff}
P_{\text{in,eff}} = \max\!\bigl(P_{\text{in}},\;
  P_{\text{rad}} - P_{\text{cp}}\bigr)
\end{equation}

\noindent This ensures $P_{\text{transport}} \geq 0$ and correctly
penalizes radiation-dominated regimes through increased recirculating
power: the additional heating must be supplied at wall-plug efficiency
$\eta_{\text{aux}}$, which raises $P_{\text{recirc}}$ and degrades the
engineering~$Q$. For example, a p-${}^{11}$B mirror at
$T_e = 150$~keV requires $P_{\text{in,eff}} \gg P_{\text{in}}$ due to
synchrotron losses, driving the LCOE well above D-T levels.

\subsubsection{Thermal and electric conversion.}

Neutron power is absorbed in the breeding blanket, where nuclear
reactions (primarily ${}^6$Li(n,$\alpha$)T) release additional energy
characterized by the blanket multiplication factor $M_n$ (typically
1.1 for lithium blankets). The total thermal power delivered to the
heat-transfer system is:

\begin{equation}
\label{eq:p-thermal}
P_{\text{th}} = M_n \, P_n + P_{\text{rad}} + P_{\text{transport}}
              + \eta_p \, P_{\text{pump}}
\end{equation}

\noindent where $P_n$ is the neutron power, $P_{\text{rad}}$ is the
radiated power reabsorbed by the first wall, $P_{\text{transport}}$ is the
charged-particle transport power deposited on plasma-facing components
(\cref{eq:plasma-balance}), and $\eta_p P_{\text{pump}}$ is the fraction
of pumping power deposited as heat. The heating power $P_{\text{in}}$ does not appear
separately; it enters the plasma and exits via the radiation and
transport channels already included in $P_{\text{th}}$.
A conventional Rankine cycle with thermal efficiency $\eta_{\text{th}}$
converts thermal power to gross electric power $P_{\text{et}}$:

\begin{equation}
\label{eq:p-gross}
P_{\text{et}} = \eta_{\text{th}} \, P_{\text{th}}
\end{equation}

\noindent For systems with direct energy conversion (DEC), a fraction
$f_{\text{dec}}$ of the transport power bypasses the
thermal cycle and is converted at efficiency $\eta_{\text{dec}}$, and
the gross electric power becomes
$P_{\text{et}} = \eta_{\text{th}} P_{\text{th}} +
 f_{\text{dec}} \, \eta_{\text{dec}} \, P_{\text{transport}}$.
DEC is disabled by default ($f_{\text{dec}} = 0$) pending further
technology maturation.

\subsection{Recirculating Power and Net Electric Output}

A fraction of the gross electric power $P_{\text{et}}$ is recirculated
to sustain the plant:

\begin{equation}
\label{eq:p-recirc}
P_{\text{recirc}} = \frac{P_{\text{in,eff}}}{\eta_{\text{aux}}}
  + P_{\text{coils}} + P_{\text{pump}} + P_{\text{cryo}}
  + P_{\text{cool}} + f_{\text{sub}} \, P_{\text{et}}
  + P_{\text{other}}
\end{equation}

\noindent where $P_{\text{in,eff}}/\eta_{\text{aux}}$ is the wall-plug
power for the heating system (using the effective heating power from
\cref{eq:p-input-eff}), $P_{\text{coils}}$ powers the magnets,
$P_{\text{cryo}}$ and $P_{\text{cool}}$ serve the cryogenic and cooling
systems, $f_{\text{sub}} P_{\text{et}}$ covers balance-of-plant
subsystems (typically 3\%), and $P_{\text{other}}$ includes tritium
processing and housekeeping loads. The net electric power is:

\begin{equation}
\label{eq:p-net}
P_{\text{net}} = P_{\text{et}} - P_{\text{recirc}}
\end{equation}

\noindent The engineering gain $Q_{\text{eng}} = P_{\text{et}} /
P_{\text{recirc}}$ must exceed unity for the plant to produce net
electricity. The recirculating fraction
$1/Q_{\text{eng}}$ is a key figure of merit: lower recirculating
fractions leave more power for sale and directly reduce the LCOE
denominator.

\subsection{Pulsed Power Balance}
\label{sec:pulsed-balance}

Pulsed confinement concepts (laser IFE, Z-pinch, and inductive
recovery schemes) differ from the steady-state balance above in two
respects: (i)~energy is delivered in discrete pulses rather than as
continuous heating, and (ii)~charged-particle energy may be recovered
electromagnetically rather than thermally. \textsc{1costingfe}
provides two pulsed power balance variants, selected by the
\texttt{PulsedConversion} parameter.

\subsubsection{Per-pulse energy framework.}
The user specifies three primary inputs: the target net electric power
$P_{\text{net}}^*$, the engineering gain $Q_{\text{eng}} = P_{\text{et}}
/ P_{\text{recirc}}$, and the repetition rate $f_{\text{rep}}$ (Hz).
From these, the model derives the required driver energy per pulse
$E_{\text{drv}}$ (MJ), the average driver power
$P_{\text{drv}} = E_{\text{drv}} \, f_{\text{rep}}$, and the fusion
power $P_{\text{fus}}$.  The scientific gain
$Q_{\text{sci}} = P_{\text{fus}} / P_{\text{drv}}$ falls out as a
derived quantity.  The driver wall-plug efficiency $\eta_{\text{pin}}$
relates stored energy to delivered energy:
$E_{\text{stored}} = E_{\text{drv}} / \eta_{\text{pin}}$. Note that
$\eta_{\text{pin}}$ characterises pulsed-driver hardware (laser, ion
beam, capacitor bank) and is distinct from the steady-state auxiliary
heating efficiency $\eta_{\text{aux}}$ used in \cref{eq:p-recirc}.
The charged-particle fraction $f_{\text{cp}}$ is derived from the fuel
type and determines the ash--neutron split, while
$f_{\text{rad}}$ is the fraction of ash power radiated during the burn.

$Q_{\text{eng}}$ is preferred over $Q_{\text{sci}}$ as the primary input
because it directly determines economic viability ($Q_{\text{eng}} \le 1$
means no net power; $Q_{\text{eng}} \approx 2$ means 50\% recirculation),
and it works uniformly across thermal and DEC conversion pathways without
requiring concept-dependent parameterization.

\subsubsection{Pulsed thermal conversion.}
For laser IFE and Z-pinch concepts the driver energy thermalises into
the blanket alongside neutron and radiation power.  The thermal pool
is:
\begin{equation}
P_{\text{th}} = M_n P_n + P_{\text{cp}} + P_{\text{drv}} + \eta_p \, P_{\text{pump}}
\end{equation}
where $M_n$ is the neutron energy multiplier and $\eta_p$ is the
pump-power thermalisation fraction (matching the steady-state
convention of \cref{eq:p-thermal}).  All electrical output
comes from the thermal cycle at efficiency $\eta_{\text{th}}$, and
the full driver draw $P_{\text{drv}}/\eta_{\text{pin}}$ enters the
recirculating load.

\subsubsection{Pulsed inductive DEC conversion.}
For inductive recovery concepts (e.g.\ pulsed FRCs with capacitor-bank-driven compression), the driver
energy circulates in an electromagnetic loop: capacitor bank $\to$
compression coils $\to$ plasma $\to$ coils $\to$ capacitor bank.
Charged-particle energy does PdV work on the expanding plasmoid,
which is recovered inductively.  The fraction of net charged-particle
power that couples as PdV work is
\begin{equation}
f_{\text{pdv}} = 1 - \left(\frac{r_{\min}}{r_{\max}}\right)^{\!2(\gamma-1)}
\end{equation}
with $\gamma = 5/3$ for an ideal monatomic plasma.  The recovered and
net DEC powers are then
\begin{align}
P_{\text{recovered}} &= \eta_{\text{dec}} \bigl(P_{\text{drv}}
  + f_{\text{pdv}} \, P_{\text{ch,net}}\bigr), \\
P_{\text{dee}} &= P_{\text{recovered}} - P_{\text{drv}},
\end{align}
where $P_{\text{ch,net}} = (1 - f_{\text{rad}})\,P_{\text{cp}}$
and $\eta_{\text{dec}}$ is the round-trip efficiency of the
inductive recovery circuit.

A key distinction from the steady-state balance is the recirculating
load: because the capacitor bank energy goes out and comes back each
cycle, only the charging \emph{losses}
$P_{\text{drv}}\,(1/\eta_{\text{pin}} - 1)$ are drawn from the grid,
not the full $P_{\text{drv}}/\eta_{\text{pin}}$.

Neutrons, bremsstrahlung, and undirected charged particles still enter
a thermal blanket.  When $\eta_{\text{th}} > 0$ a conventional thermal
cycle runs in parallel with the DEC, and the total gross electric is
$P_{\text{et}} = P_{\text{dee}} + P_{\text{the}}$.  When
$\eta_{\text{th}} = 0$ the concept is purely inductive and all
thermal loads (pumps, balance of plant (BOP)) are zeroed automatically.

For inductive DEC, the capacitor bank is costed on a \$/J stored-energy
basis (Account~22.01.07, \cref{sec:cas2207}): $C_{220107} = c_{\text{cap}}
\times E_{\text{stored}}$, where $c_{\text{cap}}$ is the all-in installed
cost per joule.  The NOAK default and its justification are discussed in
\cref{sec:cas2207}.

For the DEC inverse, the model takes the scientific gain $Q_{\text{sci}}$
as input and derives the required stored energy per pulse:
\begin{equation}
E_{\text{stored}} = \frac{P_{\text{net}}}{f_{\text{rep}} \bigl(
  \eta_{\text{eff}}\,(1 + Q_{\text{sci}}\,f_{\text{cp}}) - 1 \bigr)
  \,\eta_{\text{pin}}}
\end{equation}
where $\eta_{\text{eff}}$ combines the DEC and thermal conversion
efficiencies.  This ensures that higher DEC efficiency correctly reduces
the required capacitor bank size rather than inflating it.

\subsection{Inverse Power Balance}
\label{sec:inverse-balance}

The costing model operates in ``inverse'' mode: given a target net
electric output $P_{\text{net}}^*$, determine the required fusion power
$P_{\text{fus}}$.

\subsubsection{Pulsed inverse.}
For pulsed concepts, $Q_{\text{eng}}$ determines the gross electric and
recirculating powers directly:
\begin{align}
P_{\text{et}} &= P_{\text{net}}^* \cdot \frac{Q_{\text{eng}}}{Q_{\text{eng}} - 1}, \\
P_{\text{recirc}} &= P_{\text{et}} / Q_{\text{eng}}.
\end{align}
The driver power follows from the recirculating budget.  For thermal
conversion, $P_{\text{recirc}} = P_{\text{drv}}/\eta_{\text{pin}}
+ P_{\text{fixed}}$; for inductive DEC,
$P_{\text{recirc}} = P_{\text{drv}}(1/\eta_{\text{pin}} - 1)
+ P_{\text{fixed}}$, where $P_{\text{fixed}}$ collects
pumping, subsystem, auxiliary, cryogenic, target, and coil loads.
Then $E_{\text{drv}} = P_{\text{drv}} / f_{\text{rep}}$, and
$P_{\text{fus}}$ is obtained from the forward energy balance
with the derived $E_{\text{drv}}$.  The inversion is closed-form
(linear in $P_{\text{drv}}$ and $P_{\text{fus}}$).

\subsubsection{Steady-state inverse.}  Substituting the forward power balance into
\cref{eq:p-net} gives $P_{\text{net}}(P_{\text{fus}})$; the inverse
problem is to find the root
$P_{\text{net}}(P_{\text{fus}}) = P_{\text{net}}^*$.

For fixed plasma parameters $(n_e, T_e, B)$, the radiated power
$P_{\text{rad}}$ is independent of $P_{\text{fus}}$.  However, the
effective heating power (\cref{eq:p-input-eff}) introduces a
nonlinearity: when $P_{\text{rad}} > P_{\text{cp}} + P_{\text{in}}$,
the effective heating depends on $P_{\text{cp}} = f_{\text{cp}} \,
P_{\text{fus}}$ (where $f_{\text{cp}}$ is the charged-particle
fraction), making $P_{\text{in,eff}}$ a function of $P_{\text{fus}}$.
The net electric power is therefore piecewise linear in
$P_{\text{fus}}$, with two regimes:

\begin{enumerate}
\item \textbf{Uncapped} ($P_{\text{rad}} < P_{\text{cp}} + P_{\text{in}}$):
  $P_{\text{in,eff}} = P_{\text{in}}$ (constant); the forward balance
  is linear in $P_{\text{fus}}$ and admits a closed-form solution.
\item \textbf{Radiation-limited} ($P_{\text{rad}} \geq P_{\text{cp}} + P_{\text{in}}$):
  $P_{\text{in,eff}} = P_{\text{rad}} - f_{\text{cp}}\,P_{\text{fus}}$;
  the forward balance is linear with a different slope (increasing
  $P_{\text{fus}}$ provides more ash power to offset radiation, reducing
  required heating).
\end{enumerate}

\noindent The inverse solver uses Newton iteration with an analytical
piecewise-constant Jacobian, initialized from the closed-form uncapped
solution.  Because $P_{\text{net}}(P_{\text{fus}})$ is piecewise linear,
convergence is exact in at most two iterations.  The Newton step is:

\begin{equation}
P_{\text{fus}}^{(k+1)} = P_{\text{fus}}^{(k)}
  - \frac{P_{\text{net}}(P_{\text{fus}}^{(k)}) - P_{\text{net}}^*}
         {\mathrm{d}P_{\text{net}}/\mathrm{d}P_{\text{fus}}\big|_{P_{\text{fus}}^{(k)}}}
\end{equation}

\noindent where the derivative is computed analytically for each regime
rather than via automatic differentiation, avoiding nested
differentiation during sensitivity analysis.


%% ------------------------------------------------------------
\section{Cost Account Structure}
\label{sec:cas}

The cost account structure maps the financial framework of
\cref{sec:lcoe} onto specific plant systems and activities. Capital
accounts CAS10--CAS60 sum to the total capital investment (CAPEX);
annualized operating accounts CAS70 and CAS80 constitute OPEX; and
CAS90 annualizes the CAPEX via the CRF. \Cref{tab:cas-overview}
shows the full hierarchy.

\begin{table}[htbp]
\caption{Cost account structure overview. Accounts with dedicated
methodology subsections are marked with~$\dagger$.}
\label{tab:cas-overview}
\centering
\begin{tabular}{llp{7.0cm}}
\toprule
CAS & Description & Method \\
\midrule
10$^\dagger$ & Pre-construction & Land intensity scaling, fuel-dependent licensing \\
20 & Total direct costs & Sum of CAS21--29 \\
21$^\dagger$ & Buildings \& structures & Per-building, per-fuel, industrial-grade (18 buildings) \\
22$^\dagger$ & Reactor plant equipment & Component-level parametric (18 sub-accounts) \\
\multicolumn{3}{l}{\textit{Per-module accounts (×\,$n_{\text{mod}}$):}} \\
\hspace{1em}22.01.01 & First wall + blanket & Volume $\times$ thermal intensity, fuel-dependent \\
\hspace{1em}22.01.02 & Shield & Volume $\times$ fuel-dependent neutron scale \\
\hspace{1em}22.01.03$^\dagger$ & Coils (magnets) & Conductor kAm $\times$ \$/kAm $\times$ markup; \$0 for IFE \\
\hspace{1em}22.01.04$^\dagger$ & Heating (MFE) / driver (pulsed) & Per-MW linear; concept-specific driver capital for pulsed \\
\hspace{1em}22.01.05 & Primary structure & Volume $\times$ power scale \\
\hspace{1em}22.01.06 & Vacuum system & Volume $\times$ power scale \\
\hspace{1em}22.01.07$^\dagger$ & Power supplies & Power-scaled (MFE); \$/J stored basis for all pulsed \\
\hspace{1em}22.01.08 & Divertor / target factory & Thermal-scaled (MFE) or power-scaled (IFE) \\
\hspace{1em}22.01.09 & Direct energy converter & Mirror/FRC override; circuit-derived markups for inductive DEC \\
\hspace{1em}22.01.10$^\dagger$ & Remote handling \& maintenance & Fuel- and concept-dependent \\
\hspace{1em}22.01.11 & Installation labor & 14\% of reactor subtotal \\
\hspace{1em}22.01.12$^\dagger$ & Isotope separation plant & Zeroed; market purchase in CAS80 \\
\multicolumn{3}{l}{\textit{Plant-wide accounts:}} \\
\hspace{1em}22.02 & Main \& secondary coolant & Power-scaled \\
\hspace{1em}22.03 & Auxiliary cooling + cryoplant & Power-scaled + cryo load basis \\
\hspace{1em}22.04 & Radioactive waste management & Thermal-scaled \\
\hspace{1em}22.05 & Fuel handling \& storage & Fuel-dependent, power-scaled (0.7) \\
\hspace{1em}22.06 & Other reactor plant equip.\ & Power-scaled (0.8) \\
\hspace{1em}22.07 & Instrumentation \& control & Thermal-scaled (0.65) \\
23$^\dagger$ & Turbine plant equipment & Linear scaling on thermal electric $P_{\text{the}}$ (NETL-calibrated) \\
24$^\dagger$ & Electric plant equipment & Linear scaling on gross electric (NETL-calibrated) \\
25$^\dagger$ & Miscellaneous plant equip.\ & Linear scaling on gross electric (NETL-calibrated) \\
26$^\dagger$ & Heat rejection systems & Linear scaling on total thermal $P_{\text{th}}$ (NETL-calibrated) \\
27$^\dagger$ & Special materials & Fuel-dependent blanket fill inventory \\
28$^\dagger$ & Digital twin & Fixed cost (\$5M) \\
29$^\dagger$ & Contingency on direct costs & 10\% FOAK, 0\% NOAK (Gen-IV EMWG) \\
30$^\dagger$ & Indirect service costs & Fraction of CAS20, construction-time scaled \\
40$^\dagger$ & Owner's costs & Fuel-dependent, staffing-based, power-law scaled \\
50$^\dagger$ & Supplementary costs & Fuel-dependent sub-account model (6 components) \\
60$^\dagger$ & Interest during construction & Uniform-spending IDC formula \\
70$^\dagger$ & Annualized O\&M + replacement & Growing-annuity levelization \\
80$^\dagger$ & Annualized fuel cost & Fuel-specific consumables, burn-fraction corrected, levelized \\
90$^\dagger$ & Annualized financial costs & CRF $\times$ total capital \\
\bottomrule
\end{tabular}
\end{table}

The following subsections detail the methodology for each account
that has undergone independent review and justification.


%% ------------------------------------------------------------
\subsection{CAS10: Pre-Construction Costs}
\label{sec:cas10}

CAS10 covers land acquisition, site and plant permits, licensing,
engineering studies, and pre-construction reports. It is typically less
than 1\% of total capital; individual sub-accounts do not warrant
detailed parametric modeling.

\subsubsection{Land and land rights (C11).}
Modern compact fusion plant designs require 20--200 acres, far less
than the legacy 1{,}000-acre assumption in earlier ARIES studies.
A representative compact pilot-plant filing \citep{cfs_arc} plans 100 acres for 400~MWe,
yielding a land intensity of \refLandIntensity~acres/MWe. At US
industrial-zoned land costs of approximately
\$\pgfmathprintnumber[1000 sep={{\,}}, precision=0]{\refLandCost}/acre
(consistent with USDA farmland averages scaled for
industrial zoning\footnote{USDA NASS, \emph{Land Values 2024 Summary}, August 2024.}), land cost for a 1~GWe plant is approximately
\$\pgfmathprintnumber[precision=1]{\refLandTotal}M, negligible in the context of multi-billion-dollar capital
investments.

The compact siting is enabled by the absence of an exclusion zone:
the NRC's 2023 decision to regulate D-T fusion under 10~CFR~Part~30
(byproduct materials) rather than Part~50/52 (reactor licensing)
eliminates the large-radius emergency planning zone required for
fission plants \citep{nrc2023}. Aneutronic fuels such as p-${}^{11}$B
likely fall outside NRC jurisdiction entirely.

\subsubsection{Plant licensing (C13).}
Licensing timelines are fuel-dependent and affect first-of-a-kind
(FOAK) builds by extending the total project time, which flows into
the levelized annual cost calculation for CAS70 and CAS80.
\Cref{tab:licensing} summarizes the adopted values.

\begin{table}[htbp]
\caption{Licensing timeline by fuel cycle.}
\label{tab:licensing}
\centering
\begin{tabular}{lll}
\toprule
Fuel & Licensing time (yr) & Rationale \\
\midrule
D-T   & 2.0 & Part~30 licensing, upper end of 1--2~yr range \\
D-D   & 1.5 & Reduced tritium inventory \\
D-${}^3$He & 0.75 & Minimal radioactivity \\
p-${}^{11}$B & 0.0 & No NRC jurisdiction \\
\bottomrule
\end{tabular}
\end{table}


%% ------------------------------------------------------------
\subsection{CAS21: Buildings and Structures}
\label{sec:cas21}

CAS21 covers all buildings and structures on the fusion plant site.
Each building is priced individually per fuel type based on its
physical scope and construction grade, not a blanket multiplier.
All fusion plants are regulated under 10~CFR~Part~30 (byproduct
materials), not Part~50 (reactor license), which means no Seismic
Category~I requirements, no NQA-1 program for most buildings, and
no fission-style containment structure.

\begin{table}[htbp]
\caption{CAS21 buildings cost by fuel cycle (1~GWe reference).}
\label{tab:cas21-fuel}
\centering
\begin{tabular}{lrrl}
\toprule
Fuel & \$/kW & M\$ at 1~GWe & Construction grade \\
\midrule
D-T    & 502 & 577 & Enhanced industrial (tritium confinement, hot cell, rad-HVAC) \\
D-D    & 446 & 513 & Enhanced industrial (reduced tritium/activation scope) \\
D-${}^3$He & 339 & 390 & Light enhanced industrial \\
p-${}^{11}$B & 308 & 354 & Industrial (no hot cell, no tritium, standard HVAC) \\
\bottomrule
\end{tabular}
\end{table}

\subsubsection{Fuel-dependent scoping.}
Each building is costed for what it physically is.  Key
differentiators by fuel:

\begin{itemize}
\item \textbf{Hot cell:} D-T \$104M (full remote handling, 3--4~ft
  shielding, NQA-1 on tritium-wetted systems); D-D \$78M (3/4 scope,
  NQA-1 overhead similar); D-${}^3$He \$23M (shielded inspection room
  only); p-${}^{11}$B \textbf{\$0} (does not exist; nothing
  activates).  The hot cell is physically integrated within the
  reactor building, following UK pilot-plant cost models, not a separate
  structure.

\item \textbf{Reactor building:} D-T \$138M (biological shielding
  walls, tritium confinement barriers); p-${}^{11}$B \$81M (standard
  heavy industrial crane hall, no shielding, personnel walk-in access).

\item \textbf{Ventilation/HVAC:} D-T \$17M (rad-HVAC: HEPA, stack
  monitoring, emergency isolation); D-D \$15M (full rad-HVAC for
  tritium zones); p-${}^{11}$B \$3.5M (standard industrial HVAC).

\item \textbf{Security:} Tritium is not listed in 10~CFR Part~37
  Category~1/2 table, so no armed response requirements apply for any fusion
  fuel.  All fuels use standard industrial security (\$2--3.5M).

\item \textbf{Site improvements:} D-T \$85M (larger footprint from
  hot cell, tritium processing); p-${}^{11}$B \$69M (standard industrial site).
\end{itemize}

\subsubsection{Scaling basis.}
Different buildings scale with different physical drivers.  The
turbine building scales with thermal electric power $P_{\text{the}}$
(steam cycle only, not DEC); the reactor building and hot cell scale
with fusion power $P_{\text{fus}}$; site improvements, control room,
and security are approximately fixed; BOP electrical buildings scale
with gross electric $P_{\text{et}}$.

\subsubsection{Comparison with benchmarks.}
\Cref{tab:cas21-benchmarks} situates the fuel-dependent CAS21 figures
of \cref{tab:cas21-fuel} against non-fusion plant types.

\begin{table}[htbp]
\caption{CAS21 buildings \$/kW benchmark comparison.}
\label{tab:cas21-benchmarks}
\centering
\begin{tabular}{lrl}
\toprule
Plant type & Buildings \$/kW & Grade \\
\midrule
Gas combined-cycle (CCGT) & 150--250 & Industrial \\
Coal (NETL B12A) & 175--300 & Heavy industrial \\
Fission nuclear (Part~50) & 800--1{,}300 & Nuclear-grade \\
\bottomrule
\end{tabular}
\end{table}

The fusion p-${}^{11}$B figure of 308~\$/kW sits at roughly
1.5$\times$ CCGT, and the D-T figure of 502~\$/kW at roughly
2.5$\times$ CCGT.
The values are derived from-scratch using industrial construction
benchmarks \citep{cushman_wakefield_2025, netl2022, nrel_atb2024,
breakthrough2022}, superseding ARIES/pyFECONS values.


%% ------------------------------------------------------------
\subsection{CAS22.01.03: Superconducting Coils}
\label{sec:cas2203}

CAS22.01.03 covers the magnet system: toroidal field (TF), central
solenoid (CS), poloidal field (PF) coils, structural casing,
insulation, quench protection, cryostat interfaces, and testing.

\subsubsection{Conductor scaling law.}
Total conductor quantity in kA-m is estimated from the field, coil
radius, and a geometry-dependent path-length factor:
\begin{equation}
\label{eq:cas2203}
C_{220103} = \frac{G \cdot B_{\max} \cdot R_{\text{coil}}^2}
                   {\mu_0 \times 1000}
             \times c_{\text{kAm}}
             \times M_{\text{concept}}
\end{equation}

\noindent where $G$ is a geometry factor (tokamak: $4\pi^2$;
stellarator: $4\pi^2 \times f_{\text{path}}$; mirror: $16\pi$),
$c_{\text{kAm}}$ is the conductor cost per kA-metre, and
$M_{\text{concept}}$ is a manufacturing markup capturing winding,
insulation, quench protection, structural casing, cryostat, and testing.

\begin{table}[htbp]
\caption{CAS22.01.03 manufacturing markups and path factors by confinement concept.}
\label{tab:cas2203-markup}
\centering
\begin{tabular}{lrrrl}
\toprule
Concept & Markup & Path factor & Example (12~T, 1.85~m) & Rationale \\
\midrule
Tokamak     &  8$\times$ & 1.0 & \$516M & TF+CS+PF complexity \\
Stellarator & 12$\times$ & 2.0 & \$1{,}546M & Non-planar 3D coils \\
Mirror      &  2.5$\times$ & 1.0 & \$26M & Simple solenoids \\
Pulsed FRC  &  1.5$\times$ & 1.0 & n/a & Theta-pinch formation coils \\
Theta pinch &  1.5$\times$ & 1.0 & n/a & Compression coils \\
MagLIF      &  2.0$\times$ & 1.0 & n/a & Axial field solenoid \\
Mag.\ target &  1.5$\times$ & 1.0 & n/a & Guide-field solenoid \\
Plasma jet  &  1.5$\times$ & 1.0 & n/a & Guide-field solenoid \\
Orbitron    &  1.5$\times$ & 1.0 & n/a & Electrostatic confinement \\
Polywell    &  2.0$\times$ & 1.0 & n/a & Polyhedral magrid \\
IFE/Z-pinch & \multicolumn{3}{l}{\$0 (no confinement magnets)} & n/a \\
\bottomrule
\end{tabular}
\end{table}

\subsubsection{Conductor cost assumptions.}
The default conductor is REBCO high-temperature superconductor at
\$50/kA-m, an aggressive NOAK target.  Current market prices for
REBCO tape are \$150--300/kA-m; the \$50/kA-m assumption reflects
projected cost reductions from scaled-up manufacturing (analogous to
the solar PV learning curve).  For comparison, Nb$_3$Sn (the ITER
conductor) is \$7/kA-m at maturity, and NbTi (LHC heritage) is
similar.

This is a deliberate design choice: the model targets NOAK economics
for a mature fusion industry, not FOAK procurement costs.  Users
performing FOAK sensitivity analysis should override the conductor
cost to \$150--300/kA-m, which roughly triples the coil account.
Unlike B-11 enrichment (where the FOAK/NOAK gap reflects a missing
industrial supply chain requiring \$10{,}000/kg falling to \$75/kg), the
REBCO cost trajectory is a manufacturing learning curve for an
existing product, making the aggressive target more defensible.


%% ------------------------------------------------------------
\subsection{CAS22.01.04: Supplementary Heating / Primary Driver}
\label{sec:cas2204}

CAS22.01.04 covers different hardware depending on the confinement
family.

\subsubsection{Steady-state MFE.}
For tokamaks, stellarators, and mirrors, the account covers
vendor-purchased auxiliary heating systems.  Cost scales linearly with
installed power:
\begin{equation}
C_{220104} = \sum_i c_i \, P_i \quad [\text{M\$}]
\end{equation}
where $c_{\text{NBI}} = 7.06$, $c_{\text{ICRF}} = 4.15$,
$c_{\text{ECRH}} = 5.00$, $c_{\text{LHCD}} = 4.00$ (all M\$/MW,
2023\$).  The per-MW costs are calibrated to ITER procurement contracts,
adjusted from FOAK to NOAK.

\subsubsection{Pulsed concepts.}
For pulsed concepts, the account covers the primary driver capital, namely
the hardware that compresses or ignites the target.  This is the pulsed
analog of the magnet system (CAS22.01.03): the mechanism that provides
confinement.  Cost scales linearly with average driver power
$P_{\text{drv}} = E_{\text{drv}} \, f_{\text{rep}}$:
\begin{equation}
C_{220104} = c_{\text{drv}} \times P_{\text{drv}} \quad [\text{M\$}]
\end{equation}

\begin{table}[htbp]
\caption{CAS22.01.04 driver \$/MW coefficients by pulsed concept.}
\label{tab:cas2204-driver}
\centering
\begin{tabular}{lrl}
\toprule
Concept & \$/MW & Hardware \\
\midrule
Laser IFE   &  8 & Diode-pumped solid-state laser \\
Heavy ion   & 12 & RF linac + storage rings \\
MagLIF      &  6 & Laser preheat system \\
Mag.\ target &  3 & Pneumatic pistons, liquid metal loop \\
Plasma jet  &  4 & Plasma gun array \\
\midrule
Z-pinch, DPF, Staged Z-pinch & 0 & Driver is electrical (CAS22.01.07) \\
Pulsed FRC, Theta pinch       & 0 & Driver is coils (CAS22.01.03) \\
\bottomrule
\end{tabular}
\end{table}

Concepts whose driver is purely electrical (pulsed power / capacitor
bank) have their driver capital captured entirely in CAS22.01.07 on a
\$/J stored basis, avoiding double-counting.  Similarly, concepts
whose confinement is provided by coils (pulsed FRC, theta pinch) have
that hardware in CAS22.01.03.

%% ------------------------------------------------------------
\subsection{CAS22.01.07: Power Supplies}
\label{sec:cas2207}

CAS22.01.07 covers the electrical power supply systems that drive the
fusion core: magnet power supplies, pulsed power drivers, and associated
switchgear.  The account has two modes.

\subsubsection{Steady-state MFE.}
For tokamaks, stellarators, and mirrors, the cost scales with gross
electric output:
\begin{equation}
  C_{220107} = 80 \times (P_{\text{et}} / 1000)^{0.7} \quad [\text{M\$}]
\end{equation}
The \$80M reference at 1~GWe originates from the ARIES-CS design study
(\$70.6M in 2004\$ for a 1.3~GWe stellarator, approximately \$78/kW in
2024\$).  Cross-checks against GW-scale HVDC converter stations
(\$100--250/kW installed)\footnote{HVDC converter station vendor literature, Siemens Energy and Hitachi ABB Power Grids, 2022--2024.} and utility-scale solar inverters
(\$30--50/kW \citep{nrel_atb2024}) bracket the figure: fusion magnet supplies
operate at lower voltage and higher current than HVDC, which is the
cheaper regime for power electronics.  No direct vendor quote exists at
fusion-relevant parameters (1--10~kV, 10--100~kA, 100+~MW).
Uncertainty: $\pm$20\%.

\subsubsection{Pulsed concepts.}
For all pulsed concepts, the cost scales with stored energy per
pulse on a \$/J basis:
\begin{equation}
  C_{220107} = c_{\text{cap}} \times E_{\text{stored}} \quad [\text{M\$}]
\end{equation}
where $c_{\text{cap}} = 0.50$~\$/J is the NOAK all-in installed cost
(capacitors, switches, charging supplies, buswork).  This is an
aggressive target: present lab-scale pricing is \$20--50/J, a 40--100$\times$
gap.  The CATF IWG extension \citep{woodruff2026b} recommends
\$1.5--4/J for NOAK.  The \$0.50/J default assumes advanced-dielectric
capacitors manufactured at high volume, a commercial viability
requirement, not a current capability.

Capacitor banks require scheduled replacement (CAS72).  At $f_{\text{rep}}
= 1$~Hz and 85\% availability, the bank executes 27~million cycles per
year.  At the NOAK baseline lifetime of $10^8$ shots, replacement occurs
every 3.7~years.  The annualized replacement cost is computed as a
present-value sum of discrete replacements over the plant lifetime,
discounted at the project interest rate.

%% ------------------------------------------------------------
\subsection{CAS22.01.09: Direct Energy Converter (electrostatic)}
\label{sec:cas2209}

CAS22.01.09 covers add-on direct energy converters for linear devices
(mirrors and steady-state FRCs) where charged-particle exhaust escapes
along open field lines and can be decelerated to recover kinetic energy
as electricity.  Inductive DEC for pulsed-FRC architectures is captured in
CAS22.01.07 via $Q_{\text{sci}}$-derived markups on the driver bank
(\cref{sec:cas2207}); the present subsection covers the electrostatic
venetian-blind variant for steady-state mirrors.

The cost scales with DEC electric output:
\begin{equation}
  C_{220109} = c_{\text{dec}} \times (P_{\text{dee}} / P_{\text{dee,ref}})^{0.7}
  \quad [\text{M\$}]
\end{equation}
where $c_{\text{dec}} = 125$~M\$ is the NOAK all-in 2024 USD reference
cost at $P_{\text{dee,ref}} = 400$~MWe DEC electric output, and the
0.7~exponent reflects sub-linear scaling of vacuum/tank infrastructure
with throughput while grid modules and power conditioning scale roughly
linearly.

The reference figure is a bottom-up build-up of the DEC-specific
add-on hardware (excluding shared mirror infrastructure already
costed under CAS22.01.06 and the expander):
\begin{itemize}
  \item Tank: \$3--15~M, central \$7~M.
    Industrial-grade fabricated 304L stainless ($\$10\text{--}15$/kg
    cleaned and ASME~Section~VIII Div~1 certified for $10^{-5}$ to
    $10^{-6}$~Torr service).  Sizing for atmospheric collapse load
    on a 10--15~m diameter, 20--40~m long shell with ring stiffeners
    every 1--2~m gives 300--500 tonnes (low end) to 600--800 tonnes
    (conservative). Same dimensional class as commercial refinery
    hydrocracker reactors \citep{exxonmobil_rotterdam_hydrocracker}
    (4--5~m $\times$ 25--30~m, 1000--1500 tonnes, ASME~VIII Div~2
    at 200~bar internal) which run \$26--30~M; the VB tank holds
    only vacuum (1~atm external collapse load), so it sits well
    below that price.  ITER cryostat (in the range \$30--40/kg,
    nuclear-grade, double-walled) is overspec on certification.
  \item HV power conditioning (DC-AC chain): \$35--55~M, central \$45~M.
    The DEC outputs DC at multi-stage retarding potentials (30--100~kV),
    which a converter chain returns to grid AC.  MISO's 2024 MTEP24
    cost guide \citep{miso2024mtep24} reports VSC HVDC valve hall
    capital at \$167/kW for $\pm 250$~kV / 500~MW and \$178/kW for
    $\pm 400$~kV / 1.5~GW; the DEC needs the IGBT/DC-conversion
    portion only (no duplicate AC interconnect), placing it at
    roughly 60--70\% of full HVDC station cost, $\$100\text{--}120$/kW.
  \item Cryopumps: \$10--25~M, central \$15~M.
    $10^6$--$10^7$~L/s pumping speed.  ITER cryopumps procurement
    \citep{iter_cryopumps} (8 torus and cryostat cryopumps, EUR 21~M
    F4E investment, manufactured by Research Instruments / Alsymex)
    and industrial cryopump pricing\footnote{Edwards/Leybold STP-iXR class turbomolecular and cryogenic pump product datasheets, list pricing 2023.}
    (\$200--500k per 20{,}000~L/s) bracket the figure.
  \item Heat collection panels: \$10--20~M, central \$15~M.
    Hypervapotron-cooled tungsten panels for the 30--90~MW heat
    load deposited on grids and collectors (5--15\% of charged-particle
    power for geometrically transparent grids).  Design analog: the
    ITER NBI residual ion dump and calorimeter \citep{hemsworth2017},
    which handle 5--13~MW/m${}^2$ heat flux on hypervapotron-cooled
    Cu-CrZr / W elements at the 16~MW beam-line scale.  Public
    procurement values for the ITER beam-dump components are not
    disclosed; the dollar figure here is an order-of-magnitude
    estimate scaled by area and heat flux.  ITER divertor pricing
    ($\$60$M/GW${}_\text{th}$) is overspec, since VB grids are not
    exposed to direct fusion plasma or 14~MeV neutron flux.
  \item Grid and collector modules: \$10--20~M, central \$15~M.
    Tungsten ribbon grids, alumina HV insulators, stainless
    mounting frames.  Geometry-driven; this line uses
    \citet{hoffman1977} ($\$4$k/m${}^2$ in 1975 USD, $\$23$k/m${}^2$
    in 2024 USD via CPI) at 300--1000~m${}^2$ collector area.
    Cross-checked against the TMX-integrated 48\% demonstration
    \citep{barr_moir1983}.
  \item Misc (HV bushings, vacuum gate valve, controls, instrumentation):
    \$10--17~M, central \$12~M.
\end{itemize}

The hardware subtotal is \$78--152~M, central \$109~M.  Adding 15\%
installation (NOAK) gives \$90--175~M total, central \$125~M, which
is the value used for $c_{\text{dec}}$.  Uncertainty: $\pm 35\%$,
dominated by the HV power conditioning scale-up factor and the
plant-scale tank fabrication unit cost.

%% ------------------------------------------------------------
\subsection{CAS22.01.10: Remote Handling \& Maintenance Equipment}
\label{sec:cas2210}

CAS22.01.10 covers the robotic and remotely-operated equipment required
for in-vessel maintenance: manipulators, transfer casks, in-pipe
welding/cutting tools, and rad-hardened actuators. The building that
houses this equipment (hot cell) is costed separately in CAS21.

The cost model is:
\begin{equation}
\label{eq:cas2210}
C_{220110} = C_{\text{RH,base}}(f) \times S_{\text{concept}}
             \times \left(\frac{P_{\text{et}}}{1000}\right)^{0.5}
\end{equation}

\noindent where $C_{\text{RH,base}}(f)$ is a fuel-dependent base cost
(M\$ at 1~GWe tokamak reference) and $S_{\text{concept}}$ is a
geometry-dependent scaling factor.

\subsubsection{Fuel dependence.}
Neutron activation determines whether remote handling is needed and how
heavily equipment must be rad-hardened.
\Cref{tab:rh-fuel} lists the base costs.

\begin{table}[htbp]
\caption{Remote handling base costs by fuel cycle (1~GWe tokamak reference).}
\label{tab:rh-fuel}
\centering
\begin{tabular}{lrl}
\toprule
Fuel & Base (M\$) & Rationale \\
\midrule
D-T          & 150 & Full RH suite, rad-hardened to 14.1~MeV neutrons \\
D-D          & 100 & Same scope, reduced rad-hardening (2.45~MeV) \\
D-${}^3$He   &  30 & Lighter duty, activation still exceeds occupational limits \\
p-${}^{11}$B &  20 & Conventional maintenance equipment (no rad-hardening) \\
\bottomrule
\end{tabular}
\end{table}

\subsubsection{Concept dependence.}
Toroidal vessels (tokamak, stellarator) require articulated multi-axis
manipulators to reach components through narrow ports, while linear
geometries (mirror, FRC) permit simpler end-access tooling.
Toroidal concepts use $S_{\text{concept}} = 1.0$; linear concepts
use $S_{\text{concept}} = 0.55$.


%% ------------------------------------------------------------
\subsection{CAS22.01.12: Isotope Separation}
\label{sec:cas2212}

CAS22.01.12 is set to zero for all fuel cycles. All isotope procurement
is modeled as market purchase in CAS80 at enriched unit prices.

The rationale is that CAS22.01.12 (on-site separation plant capital) and
CAS80 (enriched market prices) are mutually exclusive models of the same
cost. If CAS80 pays \$2{,}175/kg for deuterium, that price already
includes the supplier's extraction and enrichment costs. Simultaneously
budgeting a \$300M on-site heavy-water plant charges for separation
twice. The correct model is one or the other:

\begin{itemize}
\item \textbf{Market purchase:} CAS80 at enriched \$/kg, CAS22.01.12 $= 0$.
\item \textbf{On-site plant:} CAS22.01.12 capital cost, CAS80 at raw
  feedstock prices.
\end{itemize}

Market purchase is adopted because on-site separation is not justified
at fusion-scale consumption rates. A 1~GWe D-T plant consumes
73--194~kg/yr of deuterium against a global enriched supply exceeding
80{,}000~kg/yr from heavy-water plants. Tritium is bred on-site in the
blanket and processed by the fuel handling system (CAS22.05); it is not a
market purchase. Lithium-6 enrichment capacity (1--2~t/yr
globally) could be stressed by a fleet of D-T reactors, but a
centralized enrichment facility shared across plants is a fleet-level
investment, not a per-plant CAS22 capital item.

\Cref{tab:fuel-prices} lists the retained CAS80 unit costs.

\begin{table}[htbp]
\caption{Fuel isotope unit costs (CAS80, market purchase).}
\label{tab:fuel-prices}
\centering
\begin{tabular}{lrl}
\toprule
Isotope & \$/kg & Basis \\
\midrule
Deuterium & 2{,}175 & STARFIRE (1980), inflation-adjusted \\
Li-6 (enriched) & 1{,}000 & 90\% enriched Li-6 \\
He-3 & 2{,}000{,}000 & Scarcity pricing \\
Protium & 5 & Commodity H$_2$ \\
B-11 (enriched) & 10{,}000 & FOAK estimate (B-10 tails) \\
\bottomrule
\end{tabular}
\end{table}


%% ------------------------------------------------------------
\subsection{CAS23--26: Balance of Plant Equipment}
\label{sec:cas23-26}

CAS23--26 cover conventional power-conversion and electrical infrastructure.
These accounts are \textbf{fuel-independent}: the choice of fuel
affects the reactor island (CAS22) but not the turbine hall, switchyard,
or cooling towers.

The model supports three thermal cycles, selected by the user: steam
Rankine, supercritical CO$_2$ (sCO$_2$) Brayton, and combined cycle
(gas topping plus steam bottoming).  The cycle determines the thermal
conversion efficiency $\eta_{\text{th}}$ and the per-MW coefficients for
CAS23 (turbine plant) and CAS26 (heat rejection).  CAS24 (electrical
equipment) and CAS25 (miscellaneous plant) are cycle-independent.

CAS23 (turbine plant) scales with thermal electric power
$P_{\text{the}} = \eta_{\text{th}} P_{\text{th}}$, and CAS26 (heat
rejection) scales with total thermal power $P_{\text{th}}$.  This ensures
both accounts auto-zero for pure direct energy conversion plants
($\eta_{\text{th}} = 0$) and scale correctly for hybrid DEC-plus-thermal
configurations.  CAS24 and CAS25 scale with gross electric $P_{\text{et}}$:
\begin{equation}
\label{eq:bop}
C_{23} = n_{\text{mod}} \, P_{\text{the}} \, c_{23}, \quad
C_{26} = n_{\text{mod}} \, P_{\text{th}} \, c_{26}, \quad
C_{24,25} = n_{\text{mod}} \, P_{\text{et}} \, c_{2x}
\end{equation}
\noindent where $c_{2x}$ is a per-MW coefficient (M\$/MW, 2024\$).
\Cref{tab:bop} lists the adopted values for each cycle.

\begin{table}[htbp]
\caption{Balance of plant per-MW coefficients by thermal cycle (2024\$).}
\label{tab:bop}
\centering
\begin{tabular}{llrrr}
\toprule
CAS & Scope & Rankine & sCO$_2$ & Combined \\
\midrule
23 & Turbine plant equipment & 0.198 & 0.155 & 0.235 \\
24 & Switchyard, transformers, cabling & 0.084 & 0.084 & 0.084 \\
25 & Fire protection, compressed air, HVAC & 0.051 & 0.051 & 0.051 \\
26 & Cooling towers, circulating water & 0.034 & 0.022 & 0.018 \\
\midrule
& $\eta_{\text{th}}$ & 0.40 & 0.47 & 0.53 \\
\bottomrule
\end{tabular}
\end{table}

The Rankine coefficients derive from the ARIES cost-account tradition
\citep{waganer2013}, calibrated against the NETL \textit{Cost and
Performance Baseline for Fossil Energy Plants} \citep{netl2022}.  Pre-inflation
(2019\$) values are escalated by a CPI factor of 1.22.

The sCO$_2$ Brayton efficiency of $\eta_{\text{th}} = 0.47$ reflects the
DOE/Sandia target of 45--50\% for recompression Brayton cycles at nuclear
heat source temperatures (500--700$^\circ$C)
\citep{doe_qtr2015_sco2,sandia_aec}.  The CAS23 coefficient is 22\% below
Rankine: sCO$_2$ turbomachinery is 85\% smaller by volume, but
recuperators add cost, netting 20\% overall reduction
\citep{ncepu_sco2_comparison}.

The combined cycle efficiency of $\eta_{\text{th}} = 0.53$ is derived from
NGCC performance ($>$60\%, \citealt{netl2022}), adjusted downward for fusion
heat source temperatures (600$^\circ$C vs.\ 1500$^\circ$C gas
turbine inlet).  The CAS23 coefficient is 19\% above Rankine due to dual
turbomachinery sets (topping plus bottoming) and a heat recovery steam generator (HRSG).

Cross-checks against the NREL ATB 2024 nuclear cost breakdown
\citep{nrel_atb2024} (3.9\% of total capital for energy conversion,
6.3\% for electrical equipment) confirm
that the adopted Rankine values are conservative, appropriate for NOAK fusion plants
using conventional, non-nuclear-grade BOP equipment.


%% ------------------------------------------------------------
\subsection{CAS27: Special Materials}
\label{sec:cas27}

CAS27 covers the initial inventory of non-fuel reactor materials:
breeding blanket fill (PbLi, Li, FLiBe), neutron multiplier (Be), and
other special inventory.  CAS22.01.01 covers the blanket \textit{structure};
CAS27 covers the \textit{material that fills it}.

\begin{equation}
\label{eq:cas27}
C_{27} = C_{\text{mat,base}}(f) \times \frac{P_{\text{net}}}{1000}
\end{equation}

The default assumes a PbLi blanket concept for D-T: 4{,}000 tonnes
of Pb-17Li eutectic at \$3/kg plus enriched lithium premium,
totaling \$15M at 1~GWe.  This is small relative to CAS22
(\$2.7B) because PbLi is an inexpensive commodity (mostly lead).

\textbf{Design-dependent override:} The Helium-Cooled Pebble Bed (HCPB)
blanket concept requires 300--490 tonnes of beryllium pebbles
(\$500--800/kg), costing \$180--300M.  FLiBe molten-salt blankets
cost \$25--300M depending on isotopic requirements.  These cases should
override CAS27 directly.

Non-D-T fuels have minimal or zero special materials: no breeding
blanket fill, no neutron multiplier, conventional coolants.


%% ------------------------------------------------------------
\subsection{CAS28: Digital Twin}
\label{sec:cas28}

CAS28 covers a plant-wide digital twin for operator training, predictive
maintenance, and operational optimization.  A fixed cost of \$5M is adopted,
independent of plant size, fuel type, or number of modules.  The estimate
is sourced from NtTau Digital LTD (in-house assessment for the pyFECONS
framework \citep{woodruff2026}), consistent with industrial digital twin
deployments for large thermal plants (\$2--10M).


%% ------------------------------------------------------------
\subsection{CAS29: Contingency on Direct Costs}
\label{sec:cas29}

CAS29 is a percentage-based contingency allowance on the sum of CAS21--28:
\begin{equation}
\label{eq:cas29}
C_{29} = f_{\text{cont}} \times \sum_{i=21}^{28} C_i
\end{equation}
\noindent where $f_{\text{cont}} = 10\%$ for FOAK and $0\%$ for NOAK,
following the Gen-IV EMWG convention \citep{emwg2007}.  The 10\% FOAK
contingency is conservative relative to the EMWG recommendation of 15--20\%
for advanced reactors; individual CAS22 sub-accounts already include
conservative fabrication markups.  The 0\% NOAK contingency reflects the
assumption of a mature, standardised design with resolved construction risk.


%% ------------------------------------------------------------
\subsection{CAS30: Capitalized Indirect Service Costs}
\label{sec:cas30}

CAS30 covers construction support services not attributable to specific
plant systems: field indirect costs (equipment rental, temporary
structures, consumables), construction supervision, and offsite design
services \citep{emwg2007}.

\subsubsection{Literature review.}
Three methodologies exist in the fusion costing literature:

\begin{enumerate}
\item \textbf{Schulte (1978):} A fixed 35\% of total direct cost (TDC),
  split as 15\% construction supervision, 15\% design services,
  5\% field indirect costs \citep{schulte1978}.

\item \textbf{Miller/ARIES LSA fractions:} A Level of Safety Assurance
  (LSA) dependent fraction of TDC, ranging from 21.7\% at LSA=1
  (inherently safe, minimal nuclear-grade requirements) to 29.0\% at
  LSA=4 (full nuclear-grade construction) \citep{miller1998, piet1986,
  logan1985}. The LSA concept \citep{piet1986} captures the cost
  reduction from eliminating nuclear-grade construction requirements
  where the physics permits.

\item \textbf{pyFECONS bottom-up:} A power-scaling formula
  calibrated from a staffing model for a 150~MWe reference plant
  \citep{woodruff2026}. This yields 5.8\% of TDC at 1~GWe, reflecting
  the labor component; non-labor indirect costs (equipment rental,
  temporary structures, consumables) that dominate CAS31 scope are
  not included in the staffing model.
\end{enumerate}

No real fusion construction data exists to calibrate CAS30. Private
fusion companies do not publish cost breakdowns at this granularity;
ITER is a one-of-a-kind research facility; and UK~STEP and EU-DEMO
cost studies use ARIES-lineage models. For context, efficient
series-built fission plants (Korean APR1400, Chinese Hualong~One)
achieve indirect costs of 12--20\% of TDC, while troubled FOAK
projects (US Vogtle~3\&4) exceed 35\%.

\subsubsection{Adopted model.}
A simpler aggregate model is adopted:

\begin{equation}
\label{eq:cas30}
\text{CAS30} = f_{\text{indirect}} \times \text{CAS20} \times
               \frac{T_c}{T_{\text{ref}}}
\end{equation}

\noindent where $f_{\text{indirect}} = \refFindirect$ (default), $T_c$ is the
construction time, and
$T_{\text{ref}} = \pgfmathprintnumber[precision=0]{\refTc}$~yr is a
reference duration. The construction-time scaling captures the primary driver
that varies between scenarios.

The default of 20\% is chosen as:
\begin{itemize}
\item Below Schulte's 35\% (1978 fission baseline, too high for NOAK fusion),
\item Below Miller's LSA=2 value of 23.2\% (ARIES tradition, likely
  somewhat high for NOAK),
\item Above pyFECONS's 5.8\% (labor-only staffing model),
\item Consistent with efficient NOAK nuclear construction internationally
  (15--20\% of TDC).
\end{itemize}

Sub-account precision (C31/C32/C35 splits) is not empirically
distinguishable for fusion and is not modeled.


%% ------------------------------------------------------------
\subsection{CAS40: Capitalized Owner's Costs}
\label{sec:cas40}

CAS40 covers one-time costs incurred by the plant owner during
construction to build an operating organization before commercial
operation: staff recruitment and training (C41), staff
housing/relocation (C42), and salary-related costs during the
pre-operational period (C43) \citep{emwg2007, prosser2024}.

\subsubsection{Relationship to CAS70.}
CAS40 and CAS70 share the same underlying staff pool but cover
different time periods. CAS40 is the capital cost of \emph{building}
the operations team (pre-COD); CAS70 is the annual cost of
\emph{running} it (post-COD). Both are derived from the same
fuel-specific staffing analysis (Section~\ref{sec:cas70}), applied
through different cost models. There is no double-counting.

\subsubsection{Methodology.}
The INL first-principles SFR cost study \citep{prosser2024}
estimates CAS40 sub-accounts as:
\begin{itemize}
\item C41: 1.5~years of salary for 110\% of staff (IAEA-recommended
  overhire for training attrition), plus a 25\% recruiting fee.
\item C42: \$40k relocation cost for 50\% of staff.
\item C43: Benefits at 57--59\% of salary during the training period.
\end{itemize}
Applying this methodology to the fuel-specific staffing estimates
yields total pre-operational costs that scale as roughly 3.2$\times$
annual labor-plus-benefits.

\subsubsection{Adopted model.}
\begin{equation}
\label{eq:cas40}
\text{CAS40} = C_{\text{owner}}(f) \times
               \left(\frac{P_{\text{net}}}{1000}\right)^{0.5}
\end{equation}

\noindent where $C_{\text{owner}}(f)$ is a fuel-dependent base cost
(M\$ at 1~GWe reference, 2023~USD) and the exponent 0.5 captures
staffing economy of scale. \Cref{tab:cas40-fuel} lists the adopted
values.

\begin{table}[htbp]
\caption{Capitalized owner's cost by fuel cycle (1~GWe reference).}
\label{tab:cas40-fuel}
\centering
\begin{tabular}{lrrl}
\toprule
Fuel & Staff & Base (M\$) & Key drivers \\
\midrule
D-T          & 117 & 39 & Tritium systems training, HP dept, radwaste, remote handling \\
D-D          &  94 & 31 & Reduced neutron/tritium scope \\
D-${}^3$He   &  69 & 23 & Light HP program, minor tritium \\
p-${}^{11}$B &  59 & 20 & Near-industrial, radiation safety officer (RSO) only \\
\bottomrule
\end{tabular}
\end{table}

\subsubsection{Power-law scaling.}
The exponent of 0.5 is derived from S-PRISM/INL staffing data in
\cref{sec:cas70}. Both accounts are driven by total staff count and
share the same exponent.

\subsubsection{Comparison with prior methods.}
pyFECONS uses the same LSA factor (\texttt{fac\_91}) for both CAS32
(construction supervision) and CAS40 (owner's costs), fundamentally
different scopes. This artifact of the old ARIES Account~91 mapping
produces CAS40 $\approx$ 11--15\% of TDC, implicitly assuming
fission-level staffing (500 per GWe). The staffing-based
approach adopted here yields 1--2\% of TDC for fusion, correctly reflecting the
4--8$\times$ reduction in headcount under Part~30 regulation.


%% ------------------------------------------------------------
\subsection{CAS50: Capitalized Supplementary Costs}
\label{sec:cas50}

CAS50 covers capitalized items beyond direct construction, indirect
services, and owner's costs that must be incurred before commercial
operation. CAS50 is decomposed into six sub-accounts, three of which
are fuel-dependent:

\begin{equation}
\label{eq:cas50}
\text{CAS50} = \underbrace{C_{51} + C_{53} + C_{54}}_{\text{fuel-independent}}
             + \underbrace{C_{52} + C_{55} + C_{56}}_{\text{fuel-dependent}}
             + C_{59}
\end{equation}

\subsubsection{Fuel-independent sub-accounts.}
\begin{itemize}
\item \textbf{C51 (Shipping/transportation):} 1.5\% of CAS20. World
  Nuclear Association reports transportation at 2\% of overnight
  cost for fission \citep{wna2025}; fusion avoids the extreme mass of
  reactor pressure vessels and containment structures.
\item \textbf{C53 (Taxes):} 1\% of CAS20. Sales and use tax on
  construction materials, net of exemptions typical for large energy
  projects (e.g.\ NY RPTL \S485, state manufacturing equipment exemptions).
\item \textbf{C54 (Construction insurance):} 1.5\% of (CAS20 + CAS30).
  Builder's risk premiums for large projects are typically 1--3\% of
  insured value. Fusion under Part~30 avoids Price-Anderson nuclear
  liability insurance required for Part~50 reactors.
\end{itemize}

\subsubsection{Fuel-dependent sub-accounts.}
\begin{itemize}
\item \textbf{C52 (Spare parts):} Initial on-site inventory at COD.
  D-T plants require spare blanket modules, first-wall tiles, and
  divertor cassettes due to neutron damage; p-${}^{11}$B plants need
  only conventional industrial spares. Expressed as a fraction of
  CAS22--28 (reactor plant equipment plus BOP).
\item \textbf{C55 (Startup fuel/inventory):} The fusion analogue of
  the fission ``first core'' (3\% of overnight cost,
  \$120M for a 1~GWe PWR). D-T requires 1--2~kg of tritium at
  \$30{,}000/g (\$40M); p-${}^{11}$B requires only
  commodity hydrogen and boron ($<$\$0.1M).
\item \textbf{C56 (Decommissioning provision):} The present value of
  the future decommissioning obligation, discounted over the plant
  lifetime. See below.
\end{itemize}

\noindent C59 (contingency) is applied at the same rate as other
accounts.

\subsubsection{Decommissioning cost analysis.}
Fission decommissioning costs \$750--1{,}250/kWe (OECD-NEA median
\$1{,}000/kWe) \citep{oecd_nea_decom2016}, driven by spent fuel
management, long-lived actinide disposal, and Part~50 regulatory
overhead. Fusion eliminates all three: no spent fuel exists, activated
materials decay to clearance levels within 50--100~years
\citep{iaea_fusion_decom2020}, and Part~30 decommissioning requirements
are far lighter.

Total decommissioning cost is estimated by starting from the fission
reference and removing inapplicable cost drivers (spent fuel
\$300/kWe, actinide waste \$200/kWe, Part~50 overhead
\$100/kWe, armed security \$100/kWe), then adding
fusion-specific costs (tritium decontamination for D-T). The
capitalized provision is the present value at a 3\% real discount
rate over a 40-year plant life (PV factor $\approx 0.31$).

\Cref{tab:cas50-fuel} summarizes the adopted values.

\begin{table}[htbp]
\caption{CAS50 fuel-dependent parameters (1~GWe reference, 2023~USD).}
\label{tab:cas50-fuel}
\centering
\begin{tabular}{lrrrr}
\toprule
Sub-account & D-T & D-D & D-${}^3$He & p-${}^{11}$B \\
\midrule
C52 spare parts (\% of CAS22--28) & 3\% & 2.5\% & 1.5\% & 1\% \\
C55 startup fuel (M\$) & 40 & 0.1 & 10 & 0.1 \\
C56 decom.\ provision (M\$) & 127 & 93 & 65 & 53 \\
\midrule
\textit{CAS50 total at 1~GWe (M\$)} & \textit{493} & \textit{391} & \textit{336} & \textit{293} \\
\textit{\% of CAS20} & \textit{12.3} & \textit{9.8} & \textit{8.4} & \textit{7.3} \\
\bottomrule
\end{tabular}
\end{table}

\subsubsection{Comparison with prior methods.}
The D-T CAS50 estimate of 8\% of total capital is consistent with the
Woodruff MIF example at 10\% \citep{woodruff2026}; the p-${}^{11}$B
estimate of 5\% reflects reduced spare parts, negligible fuel cost,
and lighter decommissioning.


%% ------------------------------------------------------------
\subsection{CAS60: Interest During Construction}
\label{sec:cas60}

Interest during construction (IDC) represents the cost of financing
the plant during the construction period, as introduced in
\cref{sec:capex-annual}.

\subsubsection{Formula.}
The uniform-spending IDC formula is derived in \cref{sec:capex-annual}
(\cref{eq:idc-lcoe}); CAS60 reuses it directly:
\begin{equation}
\label{eq:idc}
\text{IDC} = C \left[ \frac{(1+i)^{T_c} - 1}{i\,T_c} - 1 \right]
           = f_{\text{IDC}}(i, T_c) \times C
\end{equation}

\noindent At reference conditions ($i = \refWACC$,
$T_c = \pgfmathprintnumber[precision=0]{\refTc}$~yr),
$f_{\text{IDC}} = \pgfmathprintnumber[precision=3]{\refIDCfrac}$.

\subsubsection{Interaction with CAS90.}
CAS60 and CAS90 must use complementary conventions to avoid
double-counting construction-period financing. Two valid approaches
exist:

\begin{itemize}
\item \textbf{Option~A} (adopted): Explicit IDC in CAS60, plain CRF
  in CAS90.
  \begin{equation}
  \text{Total capital} = C + \text{IDC}, \quad
  \text{CAS90} = \text{CRF} \times \text{Total capital}
  \end{equation}

\item \textbf{Option~B}: No CAS60, effective CRF in CAS90.
  \begin{equation}
  \text{CAS90} = \text{CRF} \times (1+i)^T \times C
  \end{equation}
\end{itemize}

\noindent Option~A is adopted because it makes IDC visible as a separate
line item, uses the more realistic uniform-spending assumption (versus
lump-sum in Option~B), and is consistent with EMWG/ARIES cost reporting
conventions \citep{emwg2007}.

%% ------------------------------------------------------------
\subsection{CAS70: Annualized O\&M and Replacement}
\label{sec:cas70}

CAS70 comprises two sub-accounts:

\begin{itemize}
\item \textbf{CAS71}: Annual operations and maintenance, scaled from
  a fuel-dependent base cost per MW of net electric capacity.
\item \textbf{CAS72}: Annualized scheduled component replacement.
  Replacement costs are discounted to present value at the time of
  each replacement event and annualized via the CRF.  For pulsed
  inductive DEC concepts, CAS72 includes capacitor bank scheduled
  replacement, whose lifetime is derived from cycle count and
  repetition rate.
\end{itemize}

\subsubsection{CAS71: Staffing-based O\&M.}
The CAS71 O\&M coefficient is derived from a bottom-up staffing model
that estimates headcount and compensation by division (operations,
maintenance, administration, technical, offsite) for each fuel type,
then adds non-labor costs (maintenance materials, insurance, supplies,
regulatory fees, waste treatment, and general overhead). The staffing
model is grounded in two reference datasets: the S-PRISM sodium fast
reactor staffing estimate \citep{boardman2000} (494~staff for a
1,520~MWe fission plant) and INL's first-principles SFR cost
estimation \citep{prosser2024}, which scales the S-PRISM baseline to
different plant sizes. Conventional gas CCGT staffing data (25--35
staff for 1~GWe) anchors the lower bound.

The key differentiator between fuels is neutron-related operational
overhead. All fusion machines are regulated under 10~CFR~Part~30
(byproduct material), not Part~50 (reactor licensing), eliminating
the armed security force, licensed operator requirements, emergency
planning zones, and nuclear QA programs that drive fission plant
staffing. However, fuels with higher neutron production require
radiation protection programs, radwaste management, and
radiation-controlled maintenance, costs that scale roughly with
neutron flux. The resulting fuel-dependent coefficients are listed
in \cref{tab:cas71-staffing}.

\begin{table}[htbp]
\caption{CAS71 staffing-based O\&M coefficients by fuel cycle (1~GWe reference, 2023~USD).}
\label{tab:cas71-staffing}
\centering
\begin{tabular}{lrrrl}
\toprule
Fuel & Staff & O\&M (M\$/yr) & O\&M (k\$/MW/yr) & Key drivers \\
\midrule
D-T   & 117 & 52.0 & 52 & Tritium systems, HP dept, radwaste, remote handling \\
D-D   &  94 & 38.7 & 39 & Reduced neutron flux ($\tfrac{1}{3}$ D-T), smaller T inventory \\
D-${}^3$He &  69 & 26.0 & 26 & 5\% neutron fraction, minimal tritium \\
p-${}^{11}$B &  59 & 23.6 & 24 & Aneutronic, no tritium, RSO only \\
\bottomrule
\end{tabular}
\end{table}

\noindent All values are at a 1~GWe reference in 2023~USD. For
comparison, S-PRISM fission O\&M is \$50k/MW/yr and modern gas CCGT
is \$12k/MW/yr. The pB11 estimate (\$24k/MW/yr) sits at roughly
twice the conventional CCGT level, reflecting fusion-specific system
complexity (vacuum, magnets, plasma-facing components) rather than
regulatory burden.

Annual O\&M scales with plant size via a power law:
\begin{equation}
\label{eq:cas71}
\text{CAS71}_{\text{annual}} = C_{\text{O\&M}}(f) \times
               \left(\frac{P_{\text{net}}}{1000}\right)^{0.5}
\end{equation}

\noindent where the exponent 0.5 captures staffing economy of scale.
The INL SFR data \citep{prosser2024} shows staff scaling from
236~at 165~MWe to 1{,}040 at 3{,}108~MWe ($\alpha \approx 0.5$),
driven by the large fixed component in administration, technical, and
offsite divisions. The same exponent is used for CAS40
(\cref{sec:cas40}), which shares the same staffing basis.

\subsubsection{Levelization.}
Both sub-accounts are levelized using the procedure of \cref{sec:lac};
neglecting the growing annuity would underestimate by the percentage
given there.


%% ------------------------------------------------------------
\subsection{CAS80: Annualized Fuel Cost}
\label{sec:cas80}

CAS80 covers the annualized cost of fuel isotope consumables. The
annual fuel expenditure is computed from the fusion power, the
energy per reaction, and the unit cost of fuel isotopes:

\begin{equation}
\label{eq:cas80-base}
A_{\text{fuel}} = N_{\text{mod}} \cdot P_{\text{fus}}\,[\text{W}] \cdot
  \frac{8760 \times 3600 \cdot f_{\text{avail}} \cdot c_{\text{rxn}}}
       {Q_{\text{eff}} \cdot e_{\text{MeV}}}
\end{equation}

\noindent where $P_{\text{fus}}$ is in watts, $c_{\text{rxn}}$ is the cost per fusion reaction
(summing the unit costs of the consumed isotopes weighted by their
per-reaction masses), $Q_{\text{eff}}$ is the effective energy per
reaction in MeV (fuel-dependent; see \cref{tab:fusion-reactions}), and
$e_{\text{MeV}} = 1.602 \times 10^{-13}$~J/MeV. The factor
$8760 \times 3600$ converts one year to seconds. Fuel-specific
unit costs are listed in \cref{tab:fuel-prices}.

\subsubsection{Burn fraction and fuel recovery.}
\Cref{eq:cas80-base} implicitly assumes that all injected fuel
undergoes fusion. In practice, only a fraction $\beta_{\text{burn}}$
of the injected fuel fuses per plasma pulse or confinement cycle;
the remainder must be exhausted, processed, and either recycled or
repurchased. The fuel recovery fraction $\eta_{\text{rec}}$ captures
the efficiency of recycling unburned fuel from the exhaust stream.
The effective annual fuel cost becomes:

\begin{equation}
\label{eq:cas80-burn}
A_{\text{fuel}}^{\text{eff}} = A_{\text{fuel}} \left[
  1 + \frac{1 - \beta_{\text{burn}}}{\beta_{\text{burn}}}
      \left(1 - \eta_{\text{rec}}\right)
\right]
\end{equation}

\noindent The multiplier equals unity when either
$\beta_{\text{burn}} = 1$ (complete burn) or
$\eta_{\text{rec}} = 1$ (perfect fuel recovery). For inexpensive
fuels (deuterium at \$2{,}175/kg), the correction is negligible.
For expensive fuels (particularly He-3 at \$2{,}000{,}000/kg),
the multiplier is economically significant: at
$\beta_{\text{burn}} = 0.10$ and $\eta_{\text{rec}} = 0.95$,
the effective fuel cost is $1.45\times$ the full-burn estimate;
at $\beta_{\text{burn}} = 0.05$ and $\eta_{\text{rec}} = 0.90$,
it rises to $2.9\times$.

Default values are $\beta_{\text{burn}} = 0.05$ and
$\eta_{\text{rec}} = 0.95$, reflecting typical MCF plasma
conditions where single-pass burn fractions of 1--10\% are
expected and modern exhaust fuel processing achieves high
recovery rates.

\subsubsection{Levelization.}
CAS80 is levelized using the same growing-annuity procedure as
CAS70 (\cref{sec:lac}).


%% ------------------------------------------------------------
\subsection{CAS90: Annualized Financial Costs}
\label{sec:cas90}

CAS90 converts the total capital investment into an equivalent annual
payment:

\begin{equation}
\label{eq:cas90}
\text{CAS90} = \text{CRF}(i, n) \times \text{Total capital}
\end{equation}

\noindent where total capital includes IDC (CAS60). This uses the
plain CRF (\cref{eq:crf}), not the ``effective CRF''
$\text{CRF} \times (1+i)^{T_c}$ found in some references
\citep{miller1998}. As discussed in \cref{sec:cas60}, the effective
CRF is a valid convention only when CAS60 is zero; combining it with
an explicit CAS60 double-counts construction financing.

In a constant-dollar framework, CAS91 (operating-period escalation)
is zero by convention, and CAS92 (annual regulatory fees) is not
modeled. The framework thus equates CAS90 with CAS93 (cost of money).


%% ============================================================
\section{Conclusion}
\label{sec:conclusion}

A parsimonious differentiable framework for estimating the levelized
cost of electricity from fusion power plants has been described.
The framework spans the four candidate fuel cycles and the principal
confinement families through a common power-balance and cost-account
backbone, with each account justified independently from procurement
and benchmark data rather than from heritage scaling factors alone.
The use of JAX as the computational backend exposes exact gradients
and supports vectorised sensitivity sweeps, making LCOE-driven design
exploration tractable.

Several limitations remain. Only a single confinement-specific physics
layer (the 0D tokamak model of \cref{sec:tokamak-model}) is implemented;
analogous models for stellarators, mirrors, FRCs, and inertial systems
are planned but not yet exercised in this paper. Calibration cross-checks
against the ARC and ARIES-AT reference designs are presented in
\cref{sec:benchmark}; cross-validation against pyFECONS reference cases
at the per-account level is left to future work.
Gradient-enabled use cases (sensitivity tornados, target-driven inverse
design, autodiff-versus-finite-difference comparisons) are supported by
the framework but are not demonstrated in the present text.

\section{Code Availability}
\label{sec:code-availability}

The \textsc{1costingfe} source code, example scripts, and per-account
documentation are released at
\url{https://github.com/1cfe/1costingfe}. The framework targets
Python~3.11 or later and depends on JAX. Reproduction of the figures
and account-level cost tables in this paper requires only the
example scripts shipped in the \texttt{examples/} directory.


%% ============================================================

\appendix

\section{0D Tokamak Model}
\label{sec:tokamak-model}

The power balance of the physics module accepts fusion power
$P_{\text{fus}}$ as an input. To derive $P_{\text{fus}}$ from machine
geometry and plasma parameters for a tokamak, the framework implements a
zero-dimensional (0D) equilibrium model that chains standard scaling
laws. The model is fully JAX-differentiable, enabling gradient-based
sensitivity analysis of LCOE with respect to machine-level design
choices ($R$, $a$, $B$, $q_{95}$, etc.).

\subsection{D-T Reactivity}

The D-T fusion reactivity $\langle\sigma v\rangle$ is evaluated using
the analytic parameterization of \citet{boschhale1992}, valid for
$0.2 \leq T \leq 100$~keV:

\begin{equation}
\label{eq:sigma-v}
\langle\sigma v\rangle = C_1\,\theta\,
  \sqrt{\frac{\xi}{m_r c^2\,T^3}}\;\exp(-3\xi)
  \quad [\text{m}^3/\text{s}]
\end{equation}

\noindent where $\theta = T / (1 - T(C_2 + T(C_4 + TC_6))/(1 + T(C_3 + T(C_5 + TC_7))))$,
$\xi = (B_G^2 / 4\theta)^{1/3}$, $B_G = 34.38$~keV$^{1/2}$ is the
Gamow constant, and $C_1$--$C_7$ are the tabulated coefficients from
\citet{boschhale1992}. The reactivity peaks near $T \approx 65$~keV.

\subsection{Plasma Equilibrium}

Given major radius $R$, minor radius $a$, elongation $\kappa$,
toroidal field $B$, and edge safety factor $q_{95}$, the plasma
current follows from MHD force balance:

\begin{equation}
\label{eq:ip}
I_p = \frac{2\pi\,a^2\,\kappa\,B}{\mu_0\,R\,q_{95}}
  \quad [\text{A}]
\end{equation}

The electron density is set as a fraction $f_{\text{GW}}$ of the
Greenwald density limit \citep{iterphysics1999}:

\begin{equation}
\label{eq:greenwald}
n_{\text{GW}} = \frac{I_p\,[\text{MA}]}{\pi\,a^2}
  \quad [10^{20}\;\text{m}^{-3}],
\qquad
n_e = f_{\text{GW}}\,n_{\text{GW}}
\end{equation}

The plasma volume and first-wall area of the elongated torus are
$V = 2\pi^2 R\,a^2\,\kappa$ and $A_{\text{fw}} = 4\pi^2 R\,a\,\kappa$.

\subsection{Fusion Power}

For a 50/50 D-T plasma ($n_D = n_T = n_e/2$), the volumetric fusion
power is:

\begin{equation}
\label{eq:pfus-tokamak}
P_{\text{fus}} = \tfrac{1}{4}\,n_e^2\,
  \langle\sigma v\rangle(T_e)\;E_{\text{fus}}\;V
  \quad [\text{W}]
\end{equation}

\noindent where $E_{\text{fus}} = 17.6$~MeV per D-T reaction.
This $P_{\text{fus}}$ feeds directly into the fuel-cycle-dependent
power split (\cref{tab:fusion-reactions}) and the power balance
of \cref{eq:p-thermal}--\cref{eq:p-net}.

\subsection{Energy Confinement}

The energy confinement time is evaluated against the IPB98(y,2)
H-mode scaling \citep{iterphysics1999}:

\begin{equation}
\label{eq:ipb98y2}
\tau_{E,98} = 0.0562\;
  I_p^{0.93}\,B^{0.15}\,\bar{n}_{19}^{0.41}\,
  P_{\text{heat}}^{-0.69}\,R^{1.97}\,\epsilon^{0.58}\,
  \kappa^{0.78}\,M^{0.19}
  \quad [\text{s}]
\end{equation}

\noindent where $\bar{n}_{19}$ is the line-averaged density in units of
$10^{19}$~m$^{-3}$, $P_{\text{heat}} = P_\alpha + P_{\text{aux}}$ is
the total heating power in MW, $\epsilon = a/R$ is the inverse aspect
ratio, and $M$ is the effective ion mass in AMU. The ratio of the
actual confinement time $\tau_E = W_{\text{th}} / P_{\text{heat}}$
(where $W_{\text{th}} = 3\,n_e\,T_e\,V$ for electrons and ions at
equal temperatures) to the scaling prediction gives the H-factor,
$H = \tau_E / \tau_{E,98}$, which measures the quality of confinement
relative to the empirical database.

\subsection{Stability and Engineering Limits}

The normalized beta,

\begin{equation}
\label{eq:beta-n}
\beta_N = \frac{\beta_t\,a\,B}{I_p\,[\text{MA}]},
\qquad
\beta_t = \frac{2\mu_0\,n_e\,T_e}{B^2},
\end{equation}

\noindent must remain below the Troyon limit ($\beta_N \lesssim 3.5$)
for ideal MHD stability. The safety factor $q_{95}$ must exceed
2 to avoid kink instabilities.

Two engineering constraints are evaluated: the neutron wall loading
$q_{\text{wall}} = P_n / A_{\text{fw}}$, which drives structural
damage and blanket lifetime; and the divertor heat flux, estimated
from a simplified scrape-off-layer model with midplane power width
$\lambda_q \approx 2$~mm and flux expansion factor $f_{\text{exp}}
\approx 5$.

\subsection{Disruption Model}

Disruptions destroy confinement and can cause significant structural
damage. The model estimates a disruption frequency from the proximity
of the operating point to stability boundaries:

\begin{equation}
\label{eq:disruption-rate}
\Gamma = \sum_{j} \Gamma_0\,\exp\!\left(-s\,m_j\right)
\end{equation}

\noindent where the sum runs over the Greenwald ($m = 1 - f_{\text{GW}}$),
Troyon ($m = 1 - \beta_N/\beta_{N,\max}$), and kink
($m = 1 - q_{95,\min}/q_{95}$) channels, $\Gamma_0 = 0.1$~FPY$^{-1}$
is the baseline rate far from limits, and $s = 15$ sets the
steepness. Each disruption damages a fraction of the first-wall
component lifetime and incurs 72~hours of downtime, reducing
effective plant availability.

\subsection{Operating Modes}

The model supports two operating modes. In \emph{forward mode}, the
user specifies machine geometry, field, safety factor, Greenwald
fraction, and temperature; the model returns the full plasma state
including $P_{\text{fus}}$, $\beta_N$, $H$-factor, wall loading,
and disruption rate. In \emph{inverse mode}, the user specifies a
target $P_{\text{net}}$ and the model solves for the electron
temperature $T_e$ that produces the required fusion power via
bisection, returning both the plasma state and the complete power
table. The inverse mode enables target-driven design: given a net
power requirement, the framework determines the plasma conditions,
checks them against stability and engineering limits, and computes
the resulting LCOE in a single differentiable pass.


\section{Synchrotron radiation model details}
\label{app:synchrotron}

This appendix documents the implementation details of the synchrotron
radiation model summarized in equation~\eqref{eq:p-sync}.

\subsection{Profile assumptions and central-value conversion}

The Albajar model is a model of synchrotron radiation out of a tokamak
plasma, and it requires central (on-axis) values of electron
temperature and density, whereas \textsc{1costingfe}'s power balance
operates on volume-averaged quantities. The assumed radial profiles
follow the convention of \citet{albajar2001}:
%
\begin{equation}
\label{eq:profile-Te}
T_e(\rho) = T_{e0}\,(1 - \rho^{\beta_T})^{\alpha_T}, \qquad
n_e(\rho) = n_{e0}\,(1 - \rho^{2})^{\alpha_n}
\end{equation}
%
where $\rho = r/a$ is the normalized minor radius. For profiles
depending only on $\rho$, the toroidal volume average reduces exactly
to $\langle f \rangle = 2\int_0^1 f(\rho)\,\rho\,d\rho$ (the
$(R + a\rho\cos\theta)$ Jacobian factor integrates out). For
$\beta_T = 2$ (the default) this gives
%
\begin{equation}
\label{eq:peaking}
\langle T_e \rangle = \frac{T_{e0}}{1 + \alpha_T}, \qquad
\langle n_e \rangle = \frac{n_{e0}}{1 + \alpha_n}
\end{equation}
%
The code inverts these relations to obtain central values from the
volume-averaged inputs: $T_{e0} = \langle T_e \rangle\,(1+\alpha_T)$
and $n_{e0} = \langle n_e \rangle\,(1+\alpha_n)$. The temperature
profile shape exponent is fixed at $\beta_T = 2$ (standard parabolic
profile). Default peaking exponents are $\alpha_T = 1.0$ and
$\alpha_n = 0.5$, giving peaking factors of 2.0 and 1.5, respectively.

\subsection{Profile shape factor}

The profile shape factor $K$ in equation~\eqref{eq:p-sync} is
\citep{albajar2001}, evaluated at $\beta_T = 2$:
%
\begin{equation}
\label{eq:K-factor}
K(\alpha_n, \alpha_T) =
  (\alpha_n + 3.87\,\alpha_T + 1.46)^{-0.79}\;
  (1.98 + \alpha_T)^{1.36}\;
  2^{2.14}\;
  (2^{1.53} + 1.87\,\alpha_T - 0.16)^{-1.33}
\end{equation}

\subsection{Aspect ratio correction}

The geometric correction factor $G$ accounts for the finite aspect
ratio of the torus \citep{albajar2001}:
%
\begin{equation}
\label{eq:G-factor}
G(A) = 0.93\,\bigl[1 + 0.85\,e^{-0.82\,A}\bigr]
\end{equation}
%
where $A = R/a$ is the aspect ratio. In the large aspect-ratio limit
$G \to 0.93$.

\subsection{Application to non-tokamak geometries}

For stellarators the formula is applied directly using the
stellarator's major and minor radii. For magnetic mirrors, which have
a cylindrical geometry of length $L$ and plasma radius $a$, an
effective major radius is defined as
%
\begin{equation}
\label{eq:R-eff-mirror}
R_{\text{eff}} = \frac{L}{2\pi}
\end{equation}
%
This maps the cylinder to a torus of equivalent volume
($V = 2\pi^2 R a^2 \kappa \to \pi a^2 L$ for $\kappa = 1$). The
wall reflectivity $R_w$ is reduced for mirrors (default 0.4 vs.\ 0.6
for tokamaks) to account for radiation escaping through the open ends.

\bibliographystyle{unsrtnat}
\bibliography{references}

\end{document}
